% !TEX root = ../print.tex
% Chapter 2, opening: what is cited. The draft's section 2.0 (lines 43-118 of
% draft/spatial-hemigroup-pyramids.md): the axioms on the line and in R^d, the two
% classifications, the generator form and the notation; the definitions this blueprint takes from
% the released modules and the literature; and one INTERFACE NODE per entry of the axiom ledger
% blueprint/AXIOMS.md that a node of this blueprint spends or refers to (28 entries; A15 added in
% review round 3, batch C, decision D3). 02-line.tex
% continues the chapter with the theory on the line.
%
% The interface nodes. thm:cited-A<N> states entry A<N> as the ledger's "Statement as used": the
% mathematical sentences of the source, at the page's letter as the librarian transcribed them
% from the page image (blueprint/AXIOMS-verbatim.md, from records/pyramids/VERBATIM-*.md), narrowed
% where the spending nodes use less. Each is [A] on its one entry and has no proof, except
% thm:cited-A1, which the author regraded [T] on 2026-11-14: the shared core proves it with no
% axiom, and the node is \leanok with a proof, on Formalization/Pyramids/Cited/A1.lean; and
% thm:cited-A34, regraded [T] the same day in wave 2g, as the affine export proves it with no
% axiom (Formalization/Pyramids/Cited/A34.lean). The author
% admitted seven entries as Lean axioms on 2026-11-13 (A39, A40, A41, A42, A46, A48, A51); their
% nodes are \leanok on the declarations of Formalization/Pyramids/Cited/A<N>.lean, which are
% axioms (A42, A46, A48), the affine export's axioms restated (A39), or theorems proved from them
% (A40, A41) or from a definition (A51); since review round 3 (batch C, 2026-11-23) A42 is a
% theorem with no axiom, and A48's axiom no longer carries the finiteness of nu, proved in its
% node. The author admitted seven more on 2026-11-14 (A31, A34,
% A35, A43, A45, A49, A50); their nodes are \leanok in the same way: the affine export's theorems
% (A31, A34, A35), its axioms (A45, and A43's first clause), two axioms of this module and two
% proofs (A43), an axiom of this module (A49) and a proof with no cited fact (A50).
% blueprint/trust-boundary.txt lists the twelve axiom names reached. The other twelve (thirteen
% with A15, review round 3 batch C) are not
% \leanok: their statements are typed with sorry, as
% theorems and not as axioms, in Formalization/PyramidsSkeleton/Ch02CitedFacts.lean (phase 4), and
% they enter no Lean proof. A node that spends an entry names it on
% its status line and \uses this node (LINKAGE.md rule 14). Inside an interface node's annotation
% another entry is named by its node (\ref) or in plain text, never by \ledger (LINKAGE.md rule 7).
% The definition nodes restate definitions of the released modules and of the literature at their
% printed pages. A definition asserts nothing, so they are [T] and cite no entry.
%
% Citations corrected against the librarian's records (each also marked CHANGED where it occurs):
% V-line "section 4, p. 22" -> Lemma 4.1, pp. 22-23; iso Def. 3.1 p. 10 -> pp. 10-11; iso Lem. 3.4
% p. 12 -> pp. 12-13; iso Prop. 3.11 p. 15 -> pp. 15-16; V-cone Prop. 3.3 p. 15 -> pp. 15-16; the
% cone module's Matern kernel at scale t, cited by the draft to Prop. 6.4(3), is Prop. 5.2(2), p. 43.
%
% Stays prose, and why (one line each):
% - the chapter opening and the section leads: they say what is cited and from where;
% - the stage paragraph (V-line Thm. 5.2, p. 26): orientation; no node spends it, the chapters prove
%   what they need (lem:family-structure; ledger id A2 was dropped, VERBATIM-V-line.md A2);
% - the exponents of the symmetric self-decomposable laws (V-line Prop. 2.13, p. 13): orientation;
%   the line paper states it for orientation too, and the form (2.0) enters through Theorem 2.0;
% - the gloss "the law of cX + Y_c" after def:self-decomposable: a one-line equivalent reading;
% - the correspondence of the operator and kernel axioms after thm:cited-A30: a reading of it;
% - the filtration A_{d+1} in A_d (iso Prop. 2.6, p. 8): orientation; no proof uses it;
% - the reduction behind Theorem 2.0' (iso Lem. 3.4, 3.5): how the cited theorem is proved; no node
%   spends it (lem:reduction-keeps-stratum proves the form it needs);
% - the generator form (2.0''), the Gaussian generator, the Cin ray (V-cone Prop. 6.2, 6.4(1),
%   Lemma 3.1), and the Matern and Thorin clauses (Prop. 6.4(3)-(4)): the cone's picture for the
%   reader; no node's status line spends them (the Choquet representation, Prop. 3.3, is the node
%   thm:cited-A15 since review round 3, batch C, decision D3);
% - the list of other released statements: an index, each statement stated where it is used;
% - the paragraph on the Matern ranges: a convention;
% - the notation.

\chapter{What is cited, the axioms and the line}
\label{ch:cited}

This paper rests on five releases of the author's earlier work, and this chapter states what it
takes from them, at the printed number and page. The line paper \progpaper{V-line v0.1} states the
axioms on the line and classifies the families that satisfy them. Its cone module
\progpaper{V-cone v0.1} describes these families as a convex cone, with their generators and their
computation on the line. Its selection module \progpaper{V-selection v0.1} treats locality,
non-creation of extrema and the jet. The isotropic paper \progpaper{iso v0.1} carries the axioms
and the classification to $\RR^d$. The affine paper \progpaper{affine v0.1} supplies the frame
families of chapter~6 and the affine orbits of isotropic laws that chapter~7 uses.

The first section states the axioms, the two classifications and the generator form, with the
statements of the released modules that the later chapters spend. The second states the facts of
the literature that the proofs use: the \Levy--Khintchine theory of infinitely divisible laws,
Bernstein functions, \Polya{} frequency densities and one integral of a Bessel function. Each cited
fact is a theorem below, numbered by its entry in the axiom ledger, and a proof that uses it refers
to it. The statements that the theory on the line cites stay with that theory, in the second half
of the chapter and in Appendix~\ref{ch:line-appendix}.

\section{What is cited, the axioms and the notation}
\label{sec:cited}

\subsection{The axioms on the line}
\label{subsec:cited-line}

A scale space keeps a record of a signal at every scale $t \ge 0$, the signal itself at $t = 0$,
and an operator $\Phi_{s,t}$ takes the record at scale $s$ to the record at the coarser scale $t$.
The cascade (A6) requires that two such steps compose to one. It replaces the semigroup law
$\Phi_{t+s} = \Phi_t\Phi_s$ of the classical theory, which ties the scale to a fixed additive
parametrization, and a family with (A6) is a \emph{hemigroup}. Covariance (A8) requires that a dilation
of the signal only relabel the scales, and (ND) excludes the identity.

% draft: the axioms on the line (section 2.0), lines 47-57 of draft/spatial-hemigroup-pyramids.md.
\begin{definition}[symmetric cascade measurement families]
\label{def:cascade-measurement-family}
\lean{ScaleSpace.CascadeFamily}\leanok
A \emph{symmetric cascade measurement family} is a family of linear operators
$\Phi_{s,t} : L^1(\RR) \to L^1(\RR)$, $0 \le s \le t < \infty$, such that
\begin{itemize}
\item[(A1)] each $\Phi_{s,t}$ is a bounded operator on $L^1(\RR)$;
\item[(A2)] $\Phi_{s,t}$ commutes with every translation $f \mapsto f(\cdot - a)$, $a \in \RR$;
\item[(A3)] $\Phi_{s,t}$ commutes with the reflection $f \mapsto f(-\,\cdot)$;
\item[(A4)] $\Phi_{s,t}$ maps nonnegative functions to nonnegative functions;
\item[(A5)] $\int\Phi_{s,t}f = \int f$ for every nonnegative $f \in L^1(\RR)$;
\item[(A6)] $\Phi_{t,t} = \Id$, and $\Phi_{s,t}\Phi_{r,s} = \Phi_{r,t}$ for all $r \le s \le t$
(the cascade);
\item[(A7)] for every $f \in L^1(\RR)$ the map $(s,t) \mapsto \Phi_{s,t}f$ is continuous from
$\{0 \le s \le t\}$ to $L^1(\RR)$;
\item[(A8)] there is a family $(S_\lambda)_{\lambda > 0}$ of increasing bijections of $[0,\infty)$
with $\mathcal{D}_\lambda\Phi_{s,t} = \Phi_{S_\lambda s, S_\lambda t}\mathcal{D}_\lambda$ for all
$0 \le s \le t$ and all $\lambda > 0$, where $(\mathcal{D}_\lambda f)(x) := \lambda^{-1}f(x/\lambda)$
(scale covariance);
\item[(ND)] $\Phi_{s,t} \ne \Id$ whenever $s < t$ (nondegeneracy).
\end{itemize}
\statusT\quad\emph{(A definition, \prog{V-line v0.1}{Def.~3.1, p.~17}\statedin{line:def:cascade-family}, with the actions of
translation, reflection and dilation as printed on its p.~7. It asserts nothing and cites no
ledger entry.)}
\end{definition}

The first consequence of the axioms is that each operator is a convolution.

% draft: the sentence after the axioms (section 2.0), line 59.
% CHANGED: the citation "[V-line v0.1, section 4, p. 22]" is made precise as Lemma 4.1,
% pp. 22-23 (the symmetry clause is on p. 23; VERBATIM-V-line.md, A1), and the statement is given at
% the lemma's letter, with its hypotheses (A1), (A2), (A4), (A5) in place of all of (A1)-(A8).
\begin{theorem}[each operator is a convolution]
\label{thm:cited-A1}
\lean{Pyramids.cited_A1_repr, Pyramids.cited_A1_symm, Pyramids.cited_A1_converse}
% CHANGED (2026-11-14, the author's decision): proved from the shared core, statement and proof
% are leanok; the statement is unchanged.
\leanok
\uses{def:cascade-measurement-family}
Let $\Phi_{s,t}$, $0 \le s \le t$, be linear operators on $L^1(\RR)$ with (A1), (A2), (A4) and
(A5) of Definition~\ref{def:cascade-measurement-family}. Then for each pair $0 \le s \le t$ there
is a unique probability measure $\mu_{s,t}$ on $\RR$ with
\[
  \Phi_{s,t}f = \mu_{s,t} * f, \qquad f \in L^1(\RR);
\]
under (A3) it is symmetric, invariant under $x \mapsto -x$. Conversely, every symmetric
probability measure on $\RR$ defines, by convolution, an operator with (A1)--(A5).
% CHANGED (2026-11-14, the author's decision): regraded [T]. The shared core proves the lemma with
% no axiom, so the node is a theorem of the formal development and not an interface; ledger entry
% A1 keeps the page anchor.
% CHANGED (fidelity review, 2026-11-16, R71): the status line lists the declarations' hypotheses beside the node's (LINKAGE rule 6), and says that the symmetry declaration drops (A2), (A4) and (A5).
\statusT\quad\emph{(The line paper's representation lemma,\prog{V-line v0.1}{Lemma~4.1,
pp.~22--23}: existence and uniqueness on p.~22, the symmetry under (A3) and the converse on
p.~23. Held as [T]: the shared library proves it with no axiom, and ledger entry A1 records the
page anchor and its transcription. \ref{lem:line-axioms-kernels}(1)--(2) reads it. The node says
nothing about the families, (A6)--(A8) and (ND), which that lemma translates to the kernels
itself, nor about the continuity of the transforms of the kernels, the line paper's Lemma~4.2,
which no node cites.
\textbf{Hypotheses} (the node's, each with what its declarations require): $\Phi_{s,t}$ linear operators
on $L^1(\RR)$, (A1): same (a bounded operator, the type being that of continuous linear maps of $L^1$); (A2), (A4) and (A5): same in
\texttt{cited\_A1\_repr}, weaker in \texttt{cited\_A1\_symm}, which requires none of them, so that
declaration is the stronger statement; (A3): same, required of \texttt{cited\_A1\_symm} only, as the
symmetry clause needs; the pairs $0 \le s \le t$: same, each hypothesis quantified over all of
them; the representing measure of the symmetry clause: same, a probability measure with
$\Phi_{s,t}f = \mu * f$ almost everywhere, taken as given; $\mu$ symmetric, for the converse: same,
\texttt{IsSymmetric}, $\mu$ equal to its image under $x \mapsto -x$.)}
\end{theorem}

% CHANGED (2026-11-14, the author's decision): the node is proved; the proof of record is the
% reduction to the line paper's lemma, as the shared library checks it.
% CHANGED (the theory paper, 2026-11-16): the closing sentence on checking and cited facts shortened; no step changed.
% CHANGED (the theory paper's blueprint round, 2026-11-16, review of Section 2, STATUS): the proof
% says that the released lemma, whose proof is printed there, is proved again in the formal
% development, and the closing sentence says so in place of "with no cited fact". No step changed.
\begin{proof}
\leanok
Existence and uniqueness are \prog{V-line v0.1}{Lemma~4.1, p.~22}, with the proof printed there
and proved again in the formal development:
under (A1), (A2), (A4) and (A5) each $\Phi_{s,t}$ is the convolution with exactly one probability
measure $\mu_{s,t}$, the equation $\Phi_{s,t}f = \mu_{s,t} * f$ holding in $L^1(\RR)$, that is,
almost everywhere.

For the symmetry, write $f^-(x) := f(-x)$ and $\mu^-$ for the image of $\mu = \mu_{s,t}$
under $x \mapsto -x$. The change of variables $y \mapsto -y$ gives
$\mu^- * f = (\mu * f^-)^-$ for every $f \in L^1(\RR)$. By (A3),
$\Phi_{s,t}f^- = (\Phi_{s,t}f)^-$, so $\mu^- * f = (\Phi_{s,t}f^-)^- =
\Phi_{s,t}f = \mu * f$ for every $f$. A probability measure is determined by its convolution
operator on $L^1(\RR)$: with $g(x) := (2\pi)^{-1/2}e^{-x^2/2}$, the transform of $\mu * g$ is
$\widehat\mu\,\widehat g$, and $\widehat g$ has no zero. So $\mu^- = \mu$.

For the converse let $\mu$ be a symmetric probability measure on $\RR$. The map
$f \mapsto \mu * f$ is linear, and $\|\mu * f\|_1 \le \|f\|_1$ by Tonelli's theorem, so it is a
bounded operator on $L^1(\RR)$, which is (A1). Since
$(\mu * f(\cdot - a))(x) = \int f(x - a - y)\,\mu(dy) = (\mu * f)(x - a)$, it commutes with
translation, which is (A2). Since $\mu^- = \mu$, the identity of the previous paragraph reads
$\mu * f^- = (\mu * f)^-$, which is (A3). If $f \ge 0$ almost everywhere, then
$\mu * f \ge 0$ almost everywhere, which is (A4), and Fubini's theorem gives
$\int \mu * f = \mu(\RR)\int f = \int f$, which is (A5).

Every step is machine-checked, the released lemma's included, and no fact of the literature is
cited.
\end{proof}

Under (A1)--(A7) the transform of the kernel is $\widehat\mu_{s,t} = e^{-g_{s,t}}$ with
\[
  g_{s,t}(\omega) = a(s,t)\,\omega^2 + \int_{(0,\infty)}(1 - \cos\omega x)\,\varpi_{s,t}(dx),
  \quad a(s,t) \ge 0, \quad \int_{(0,\infty)}\min(1, x^2)\,\varpi_{s,t}(dx) < \infty
\]
\prog{V-line v0.1}{Thm.~5.2, p.~26}. So every increment, the \emph{stage} $\mu_{s,t}$, is a
symmetric \emph{infinitely divisible} law: for every $n$ it is the $n$-fold convolution power of
some law. Its exponent has a Gaussian part $a(s,t)\omega^2$ and a jump part, and $\varpi_{s,t}$,
its \emph{\Levy{} measure} folded onto $(0,\infty)$, counts the jumps by their length. The pair
$(a, \varpi)$ of such an exponent is determined by the function.

% draft: none as a statement. The draft cites "[V-line v0.1, (2.2)]" and "[Prop. 2.3(3)]" in the
% proof of Lemma 2.3 (line 140), for the identification of the data from the exponent. Ledger
% entries A8 and A9 of the plan are one statement on one page (VERBATIM-V-line.md, A8, A9), so the
% display (2.2) is folded into this node and A8 is not an entry.
\begin{theorem}[the folded \Levy--Khintchine pair is unique]
\label{thm:cited-A9}
\lean{Pyramids.cited_A9}
% CHANGED (wave 3d, 2026-11-14): proved from A39(ii)-(iii), through the affine module's
% line_levy_unique; statement and the declaration's type unchanged. The entry is not admitted and
% needs no axiom of its own.
\leanok
Let $a, a' \ge 0$, and let $\nu$ and $\nu'$ be measures on $(0,\infty)$ with
$\int_{(0,\infty)}\min(1, x^2)\,\nu(dx) < \infty$ and
$\int_{(0,\infty)}\min(1, x^2)\,\nu'(dx) < \infty$. If
\[
  a\,\omega^2 + \int_{(0,\infty)}(1 - \cos\omega x)\,\nu(dx)
  = a'\,\omega^2 + \int_{(0,\infty)}(1 - \cos\omega x)\,\nu'(dx)
  \qquad\text{for every } \omega \in \RR,
\]
then $a = a'$ and $\nu = \nu'$.
% CHANGED (fourth blind pass, 2026-11-15): the derivation from A39 moves out of the annotation into
% the proof below, which names no declaration, and the statement's \uses{thm:cited-A39} moves with
% it to the proof; the Assignment lists the spenders as the \uses graph
% has them (thm:cited-A4 added, prop:discrete-gaussian-pyramid dropped, being [T] and spending
% nothing), and the sentence on the existence half says where the representation of a stage comes
% from.
% CHANGED (fidelity review, 2026-11-16, R71): the status line lists the declaration's hypotheses beside the node's (LINKAGE rule 6).
\statusA\quad\ledger{A9}\quad\emph{(Cited from \prog{V-line v0.1}{Prop.~2.3(3), p.~9}, the
symmetric \Levy--Khintchine clause with its display (2.2): every function of the line paper's
class of symmetric negative definite functions has the representation
$\psi(\omega) = a\omega^2 + \int_{(0,\infty)}(1 - \cos\omega x)\,\nu(dx)$ with $a \ge 0$ and
$\int(1 \wedge x^2)\,\nu(dx) < \infty$, ``the pair $(a,\nu)$ being unique''. This citation is the
node's grounding. Not admitted, and found to need no axiom of its own: the proof below derives the
statement from clauses (ii) and (iii) of Theorem~\ref{thm:cited-A39} read at $d = 1$.
\textbf{Assignment.} \ledger{A9} carries the statement,
which is that uniqueness clause read for two pairs defining the same function.
\ref{thm:lattice-line}(3) spends it to identify the data of the relabelled kernels, and
Theorem~\ref{thm:cited-A4} in its ``only if'', where the uniqueness of the folded pair is the one cited
input of the formal derivation. Not carried: the existence half of the clause and its converse. No
node cites them from this page. The representation of each stage is part of what
Theorem~\ref{thm:cited-A4} carries, and the formal development proves it with no cited input; a law with
given data comes from clause (iii) of A39, through \ref{lem:line-exponent-data}(1). Not carried
either: any two-sided form, the folded measure being the image of a symmetric \Levy{} measure
under $x \mapsto |x|$ by the line paper's convention, p.~10. The same uniqueness in $\RR^d$, for
laws, is Theorem~\ref{thm:cited-A39}(ii).
\textbf{Hypotheses} (the node's, each with what its declaration requires): $a, a' \ge 0$: same;
$\nu$, $\nu'$ measures on $(0,\infty)$: same (measures on $\RR$ with no mass on $(-\infty, 0]$);
$\int\min(1, x^2)\,d\nu < \infty$ and the same for $\nu'$: same (the lower integral finite); the
equality of the two displayed functions for every $\omega \in \RR$: same, of the real-valued
exponents, each finite, so it is not an equality of two junk values.)}
\end{theorem}

% CHANGED (fourth blind pass, 2026-11-15): new proof environment, the derivation formerly sketched
% in the annotation, with the symmetrization written out.
% CHANGED (the theory paper, 2026-11-16): the closing sentence on checking and cited facts shortened; no step changed.
% CHANGED (the theory paper's blueprint round, 2026-11-16, review of Section 2, PROSE): the lift to
% R^1 is no longer named in the proof; the identification of R with R^1 is stated once, in the lead
% of the theory on the line. No step changed.
\begin{proof}
\leanok
\uses{thm:cited-A39}
For a pair $(a, \nu)$ as in the statement let $\tilde\nu$ be the symmetric measure on $\RR$ with
$\tilde\nu(\{0\}) = 0$ and $\tilde\nu(B) = \tfrac12\nu(B \cap (0,\infty)) +
\tfrac12\nu(-B \cap (0,\infty))$. Then $\int\min(x^2, 1)\,\tilde\nu(dx) =
\int_{(0,\infty)}\min(x^2, 1)\,\nu(dx) < \infty$, so $(2a, \tilde\nu, 0)$ is a generating triplet
on $\RR$. In its \Levy--Khintchine bracket the imaginary part of the integrand,
$\sin\omega x - \omega x\mathbf 1_{|x| \le 1}$, is odd and $\tilde\nu$-integrable, so it integrates
to $0$. The real part $\cos\omega x - 1$ is even, so its $\tilde\nu$-integral is
$-\int_{(0,\infty)}(1 - \cos\omega x)\,\nu(dx)$. With the Gaussian term
$-\tfrac12\cdot 2a\,\omega^2$ the bracket is $-\psi(\omega)$, where $\psi$ is the function of the
pair, $\psi(\omega) = a\omega^2 + \int_{(0,\infty)}(1 - \cos\omega x)\,\nu(dx)$.

By Theorem~\ref{thm:cited-A39}(iii) there is an infinitely divisible law $\mu$ on $\RR$ with the triplet
$(2a, \tilde\nu, 0)$, so $\widehat\mu = e^{-\psi}$. By hypothesis the pair $(a', \nu')$ has the
same function $\psi$. So by the first paragraph the triplet $(2a', \tilde\nu', 0)$ also represents
$\widehat\mu$ by the formula of Theorem~\ref{thm:cited-A39}(i). By Theorem~\ref{thm:cited-A39}(ii) the two triplets
are equal: $a = a'$ and $\tilde\nu = \tilde\nu'$. On $(0,\infty)$ the measure $\tilde\nu$ is
$\tfrac12\nu$, so $\nu = 2\tilde\nu|_{(0,\infty)} = 2\tilde\nu'|_{(0,\infty)} = \nu'$.

Every step is machine-checked, and the cited fact is \ledger{A39}(ii)--(iii), at $d = 1$.
\end{proof}

The class of laws that the full axioms produce is a classical one.

% draft: the definition of self-decomposability (section 2.0), line 61; and the semi-
% selfdecomposable laws of Definition 3.4 (line 402), whose source is the same definition of Sato's.
% CHANGED (the theory paper's blueprint round, 2026-11-16, review of Section 2, POINTER): the
% residual is written varrho_b, Sato's rho_b renamed, so that the notion has one residual symbol
% with def:semi-selfdecomposable-line and def:semi-selfdecomposable. A rename.
\begin{definition}[self-decomposable and semi-selfdecomposable laws]
\label{def:self-decomposable}
\lean{ScaleSpace.IsSelfDecomposable, AffineHemigroup.IsSelfDecomposableSato,
Pyramids.IsSDResidual, Pyramids.IsSemiSelfDecomposableWithSpan, Pyramids.IsSemiSelfDecomposable}\leanok
Let $\mu$ be a probability measure on $\RR^d$. It is \emph{self-decomposable} if, for every
$b > 1$, there is a probability measure $\varrho_b$ on $\RR^d$ such that
\[
  \widehat\mu(z) = \widehat\mu(b^{-1}z)\,\widehat\varrho_b(z), \qquad z \in \RR^d.
\]
It is \emph{semi-selfdecomposable} if there are some $b > 1$ and some infinitely divisible
probability measure $\varrho_b$ with the same identity. Then $b$ is called a \emph{span} of $\mu$, and
$\varrho_b$ its \emph{residual}.
\statusT\quad\emph{(A definition, [\sscite{sato1999levy}, Def.~15.1, p.~90], formula
(15.1)\statedin{affine/iso:def:isotropic-kernel-family}; the
word ``residual'' is this paper's. It asserts nothing and cites no ledger entry. That a
self-decomposable law is infinitely divisible, with infinitely divisible residuals, is
Theorem~\ref{thm:cited-A43}.)}
\end{definition}

Writing $c = b^{-1} \in (0,1)$, a law $\mu$ is self-decomposable if for every $c$ there is a law
$\varrho_c$ with $\widehat\mu(\omega) = \widehat\mu(c\omega)\,\widehat\varrho_c(\omega)$, that is if $\mu$
is the law of $cX + Y_c$ with $X \sim \mu$ and $Y_c$ independent of $X$. The exponents of the
symmetric self-decomposable laws on the line are those of the form
\begin{equation}
\label{eq:admissible-exponent}
\tag{2.0}
  F(\omega) = a\,\omega^2 + \int_0^\infty(1 - \cos\omega x)\,\frac{k(x)}{x}\,dx
\end{equation}
with $a \ge 0$, $k : (0,\infty) \to [0,\infty)$ nonincreasing and
$\int_0^\infty\min(1, x^2)\,k(x)\,x^{-1}\,dx < \infty$ \prog{V-line v0.1}{Prop.~2.13, p.~13}, the
integrability being that of a \Levy{} measure. The classification of the line is the following
theorem.

% draft: Theorem 2.0 (section 2.0), line 65 of draft/spatial-hemigroup-pyramids.md. This node IS
% the draft's Theorem 2.0.
% CHANGED: the form (2.0) is written out in the statement, so that the node reads on its own.
\begin{theorem}[the classification on the line]
\label{thm:cited-A4}
\lean{Pyramids.cited_A4, Pyramids.cited_A4_exists, Pyramids.cited_A4_unique}
% CHANGED (wave 3d, 2026-11-14): proved from A39(ii)-(iii); the three declarations unchanged. The
% entry is not admitted and needs no axiom of its own.
\leanok
\uses{def:cascade-measurement-family, thm:cited-A9, thm:cited-A39, lem:line-exponent-data}
A family $(\Phi_{s,t})_{0 \le s \le t}$ of linear operators on $L^1(\RR)$ satisfies (A1)--(A8)
and (ND) of Definition~\ref{def:cascade-measurement-family} if and only if there are an increasing
bijection $\chi$ of $[0,\infty)$, the \emph{gauge}, and a function $F \not\equiv 0$ of the form
\[
  F(\omega) = a\,\omega^2 + \int_0^\infty(1 - \cos\omega x)\,\frac{k(x)}{x}\,dx,
\]
with $a \ge 0$, $k \ge 0$ nonincreasing and
$\int_0^\infty\min(1, x^2)\,k(x)\,x^{-1}\,dx < \infty$, such that $\Phi_{s,t}f = \mu_{s,t} * f$ with
\[
  \widehat\mu_{s,t}(\omega) = \exp\bigl(-[F(\chi(t)\omega) - F(\chi(s)\omega)]\bigr),
  \qquad 0 \le s \le t,\ \omega \in \RR.
\]
% CHANGED (skeleton review, M1): the construction clause, which the display presupposes when it
% names the measures by their transforms and which the page's following text calls the
% "construction" clause, is written out (VERBATIM-V-line.md, A4; AXIOMS-verbatim.md, A4).
Conversely, for every such pair $(\chi, F)$ there are probability measures $\mu_{s,t}$ with
these transforms, and the operators $f \mapsto \mu_{s,t} * f$ satisfy (A1)--(A8) and (ND).
The pair $(\chi, F)$ is unique once $\chi(1) = 1$.
% CHANGED (fourth blind pass, 2026-11-15): the annotation says that the printed grounding is the
% cited theorem, and describes the formal derivation by what it proves and what it cites, without
% naming declarations; the "one cited input" sentence now says that the representation of each
% stage is proved inside the shared core's analysis, which is why the uniqueness is its only input.
% CHANGED (fidelity review, 2026-11-16, R71): the status line lists the declarations' hypotheses beside the node's (LINKAGE rule 6).
\statusA\quad\ledger{A4}\quad\emph{(Cited from \prog{V-line v0.1}{Thm.~7.3, p.~37}, the main
theorem, whose $F$ is ``of the form (2.3)'' of the line paper, p.~13, with the integrability of a
\Levy{} measure: the form above. The construction sentence is the theorem's ``if'' read for the
measures its display names: the page's following text calls the three clauses ``construction,
analysis, uniqueness'' and says that what the axioms leave free is ``exactly the pair $(a,k)$''.
This citation is the node's printed grounding: the proof is the line paper's and is not reprinted
here. Not admitted, and found to need no axiom of its own: the formal development
derives the three clauses from clauses (ii) and (iii) of Theorem~\ref{thm:cited-A39}, through the shared
library's proved classification on the line. Its ``only if'' is the shared library's proof of the
line paper's analysis direction. That proof establishes the convolution form, the infinite
divisibility and the \Levy--Khintchine representation of every stage itself, by a null-array and
tightness argument with no cited input, so its one cited input is the uniqueness of the folded
pair, Theorem~\ref{thm:cited-A9}, which Theorem~\ref{thm:cited-A39}(ii)--(iii) give. The uniqueness of the
normalized pair needs no cited input once the kernels of one operator family are identified, a
measure being determined by its convolution operator. The construction and the ``if'' follow the shared library's construction,
which uses a cited fact only to obtain a law with a given symmetric exponent, and take that law from
clause (iii) (Lemma~\ref{lem:line-exponent-data}(1)) in place of Bochner's theorem.
\textbf{Assignment.} \ledger{A4} carries the whole statement: both directions and the uniqueness of
the normalized pair. \ref{thm:lattice-line}(3) spends the ``only if'',
\ref{thm:annulus-parametrization-phase}(2)--(3) the admissibility of $F$ and the form
$F(\chi(t)\,\cdot)$, and \ref{thm:matern-pyramid-line}(1) the ``if'', that an admissible
$F \not\equiv 0$ has a similarity family with the axioms. Not carried: the sentence of the page
that in the canonical gauge the kernel at scale $t$ has transform $e^{-F(t\,\cdot)}$, which is the
case $\chi = \mathrm{id}$ of the display; the monotonicity in $t$ of the data of the stages, which
the line paper's Theorem~5.2 does not state and \ref{lem:family-structure} proves; and the
identification of the folded \Levy{} measure of a kernel, which is Theorem~\ref{thm:cited-A9} with a change
of variables.
\textbf{Hypotheses} (the node's, each with what its declarations require): $\Phi$ a family of linear
operators on $L^1(\RR)$ with (A1)--(A8) and (ND): same for \texttt{cited\_A4}, whose
\texttt{CascadeFamily} requires (A1)--(A3), (A5)--(A7), positivity, nondegeneracy and (A8) at every
$\lambda > 0$; weaker for \texttt{cited\_A4\_unique}, which requires only that $\Phi$ have the kernel
families $\mu_1$, $\mu_2$ and not that it satisfy the axioms; $\chi$ a gauge: same (an increasing
bijection of $[0,\infty)$ with $\chi(0) = 0$); $F \not\equiv 0$: same for the first pair, weaker for
the second pair of \texttt{cited\_A4\_unique}, where it is not required; $a \ge 0$, $k \ge 0$
nonincreasing and the two integrability conditions: same (\texttt{SDProfile}), with one more field,
$k(0) = 0$, a normalisation that the node leaves free as $k$ is read on $(0,\infty)$, so stronger
where the profile is given and harmless where it is obtained; the kernels $\mu_{s,t}$ with the
displayed transforms: same, probability measures with $\Phi_{s,t}f = \mu_{s,t} * f$ almost
everywhere; $\chi(1) = 1$: same, for both pairs.)}
\end{theorem}

In the \emph{canonical gauge} $\chi = \mathrm{id}$ the kernel at scale $t$ has transform
$e^{-F(t\omega)}$. The functions $F$ of the form (\ref{eq:admissible-exponent}) are the
\emph{admissible exponents}, and $(a, k)$ their \emph{data}. They form a convex cone
\prog{V-cone v0.1}{Lemma~2.7, p.~12}.

\subsection{The axioms in $\RR^d$}
\label{subsec:cited-d}

Fix $d \ge 1$. For a measure $\mu$ on $\RR^d$ write
$\widehat\mu(\omega) = \int e^{i\langle\omega, x\rangle}\,\mu(dx)$, and for a linear map $A$ of
$\RR^d$ write $A_*\mu$ for the image of $\mu$ under $x \mapsto Ax$.

% draft: the axioms in R^d (section 2.0), lines 69-75; with the operator form of line 77.
\begin{definition}[isotropic cascade families]
\label{def:isotropic-cascade-family}
\lean{AffineHemigroup.IsIsotropicCascadeFamily}\leanok
Let $d \ge 1$. An \emph{isotropic cascade family in dimension $d$} is a family
$(\mu_{s,t})_{0 \le s \le t}$ of probability measures on $\RR^d$ such that
\begin{itemize}
\item[(I1)] $\mu_{t,t} = \delta_0$, and $\mu_{r,t} = \mu_{s,t} * \mu_{r,s}$ whenever
$r \le s \le t$ (cascade);
\item[(I2)] for every $\omega \in \RR^d$ the map $(s,t) \mapsto \widehat\mu_{s,t}(\omega)$ is
continuous on $\{0 \le s \le t\}$ (continuity);
\item[(I3)] $Q_*\mu_{s,t} = \mu_{s,t}$ for every $Q \in \Orth(d)$ and every pair (rotation
invariance);
\item[(I4)] there is a family $(S_\lambda)_{\lambda > 0}$ of increasing bijections of
$[0,\infty)$ with $(\lambda I)_*\mu_{s,t} = \mu_{S_\lambda s, S_\lambda t}$ for every $\lambda > 0$
and every pair (dilation covariance);
\item[(ND)] $\mu_{s,t} \ne \delta_0$ whenever $s < t$ (nondegeneracy).
\end{itemize}
\statusT\quad\emph{(A definition,
\prog{iso v0.1}{Def.~3.2, p.~11}\statedin{affine/iso:def:isotropic-kernel-family}. The isotropic paper also
states the axioms on operators on $L^1(\RR^d)$, as (A1)--(A8) and (ND) with the orthogonal group in
place of the reflection, \prog{iso v0.1}{Def.~3.1, pp.~10--11}, and joins the two forms by its
Lemma~3.3, which is Theorem~\ref{thm:cited-A30}. This blueprint states the axioms in $\RR^d$ on the kernels
only. It asserts nothing and cites no ledger entry. The declaration takes the family
$(S_\lambda)$ of (I4) as an argument, \texttt{IsIsotropicCascadeFamily $\mu$ S}; the definition's
object is $\exists S$, and every statement that uses it quantifies so. The declaration does not require
$d \ge 1$; at $d = 0$ (ND) cannot hold.)}
% CHANGED (fidelity review, 2026-11-15, R2): the status line says that the declaration takes the
% gauge S as an argument and that the definition's object is the existential.
\end{definition}

The axioms are stated on the kernels. The isotropic paper also states them on operators on
$L^1(\RR^d)$, as (A1)--(A8) and (ND) with the rotations in place of the reflection
\prog{iso v0.1}{Def.~3.1, pp.~10--11}. The two definitions describe the same families, by the
next theorem.

% draft: the sentences after the axioms in R^d (section 2.0), line 77.
% CHANGED: the citation of iso Def. 3.1 reads pp. 10-11, not p. 10 (VERBATIM-selection-iso-affine.md).
\begin{theorem}[operators and kernels in dimension $d$]
\label{thm:cited-A30}
\lean{Pyramids.cited_A30_repr, Pyramids.cited_A30_converse, Pyramids.cited_A30_continuity}
\notready
Let $d \ge 1$, and for a probability measure $\mu$ on $\RR^d$ and $f \in L^1(\RR^d)$ write
$(\mu * f)(x) := \int f(x - y)\,\mu(dy)$.
\begin{enumerate}
\item Let $\Phi$ be a linear operator on $L^1(\RR^d)$ that is bounded, commutes with every
translation $f \mapsto f(\cdot - a)$, $a \in \RR^d$, maps nonnegative functions to nonnegative
functions, and has $\int\Phi f = \int f$ for every nonnegative $f$. Then there is a unique
probability measure $\mu$ on $\RR^d$ with $\Phi f = \mu * f$ for every $f \in L^1(\RR^d)$.
Conversely, for every probability measure $\mu$ on $\RR^d$ the operator $f \mapsto \mu * f$ has
these four properties.
\item Let $\mu_n$ and $\mu$ be probability measures on $\RR^d$. Then
$\widehat\mu_n(\xi) \to \widehat\mu(\xi)$ for every $\xi \in \RR^d$ if and only if
$\mu_n * f \to \mu * f$ in $L^1(\RR^d)$ for every $f \in L^1(\RR^d)$.
\end{enumerate}
\statusA\quad\ledger{A30}\quad\emph{(Cited from \prog{iso v0.1}{Lem.~3.3(1)--(2), p.~11}; its
clause (2) is proved in the isotropic paper's supplement, \S S3.1. Its statement is typed in the skeleton
(\texttt{PyramidsSkeleton.Ch02CitedFacts}) and is not an axiom; by the author's decision of
2026-11-15 it is not admitted. \textbf{Assignment.} \ledger{A30} carries the two clauses, for
every $d \ge 1$.
% CHANGED (wave 3b, 2026-11-15): prop:line-and-arrays(1) no longer spends A30; it is proved
% through checked nodes. The entry has no spender.
No node spends it: \ref{prop:line-and-arrays}(1), which read it at $d = 1$, is proved through
the line's own bridges between operators and kernels. Not
carried: the clauses (3)--(4) of the lemma on p.~12, which no node reads.)}
\end{theorem}

So (I1) is the cascade (A6), (I2) the continuity (A7), (I3) the rotation invariance (A3) and (I4)
the covariance (A8), with the same bijections. Let $\mathcal{A}_d$ be the set of admissible $F$ for
which $\omega \mapsto F(|\omega|)$, $\omega \in \RR^d$, is the exponent of a self-decomposable law
on $\RR^d$. Then $\mathcal{A}_1$ is the whole cone of admissible exponents and
$\mathcal{A}_{d+1} \subseteq \mathcal{A}_d$ \prog{iso v0.1}{Prop.~2.6, p.~8}. The classification
in $\RR^d$ is the following theorem.

% draft: Theorem 2.0' (section 2.0), line 79 of draft/spatial-hemigroup-pyramids.md. This node IS
% the draft's Theorem 2.0'.
% CHANGED: the class A_d and the word "gauge" are written out in the statement, so that the node
% reads on its own (the isotropic paper's page refers to its section 2.2 for both).
\begin{theorem}[the classification in $\RR^d$]
\label{thm:cited-A31}
\lean{Pyramids.cited_A31, Pyramids.cited_A31_unique}
\leanok
\uses{def:isotropic-cascade-family, def:self-decomposable, thm:cited-A39}
Let $d \ge 1$, and let $\mathcal{A}_d$ be the set of admissible exponents $F$ (of the form
(\ref{eq:admissible-exponent})) for which $\omega \mapsto F(|\omega|)$ is the exponent of a
self-decomposable law on $\RR^d$. A family $(\mu_{s,t})_{0 \le s \le t}$ of probability measures
on $\RR^d$ is an isotropic cascade family in dimension $d$ if and only if there are a gauge
$\chi$, an increasing bijection of $[0,\infty)$, and an $F \in \mathcal{A}_d$ with
$F \not\equiv 0$ such that
\[
  \widehat\mu_{s,t}(\omega) = \exp\bigl(-[F(\chi(t)|\omega|) - F(\chi(s)|\omega|)]\bigr),
  \qquad 0 \le s \le t,\ \omega \in \RR^d.
\]
% CHANGED (fidelity review, 2026-11-16, R72): the sentence names the class of the second pair, $F' \in \mathcal{A}_d$, as the declaration requires; the page's wider class is in the status line.
The pair is unique against any second such pair $(\chi', F')$ with $F' \in \mathcal{A}_d$, up to
$(\chi, F) \mapsto (c\chi, F(c^{-1}\,\cdot))$ with $c > 0$, so it is unique once $\chi(1) = 1$.
% CHANGED (fidelity review, 2026-11-16, R71): the status line lists the declarations' hypotheses beside the node's (LINKAGE rule 6).
\statusA\quad\ledger{A31}\quad\emph{(Cited from \prog{iso v0.1}{Thm.~3.6, p.~13}, with the gauge
and the class $\mathcal{A}_d$ as the isotropic paper defines them in its \S 2.2, p.~6 (a gauge is
an increasing bijection $\chi$ of $[0,\infty)$, hence continuous, with $\chi(0) = 0$), and
Prop.~2.6, p.~8. Admitted as a Lean axiom by the author on 2026-11-14, and found to need none of
its own: the affine module's formal development proves the theorem
(\texttt{AffineHemigroup.isotropic\_classification} and
\texttt{isotropic\_classification\_unique}), the equivalence from clauses (ii) and (iii) of
Theorem~\ref{thm:cited-A39} and the uniqueness with no cited fact; the declarations named above are those
theorems, the uniqueness read for $F' \in \mathcal{A}_d$, which has $F'(0) = 0$. \textbf{Assignment.}
% CHANGED (wave 2g, 2026-11-14): what each declaration reaches, as printed; the node stays [A].
\ledger{A31} carries the whole statement. Of the two declarations, the uniqueness
(\texttt{cited\_A31\_unique}) rests on no axiom, and the equivalence (\texttt{cited\_A31}) reaches
clauses (ii) and (iii) of Theorem~\ref{thm:cited-A39}; the node stays [A] for the equivalence. \ref{thm:lattice-strata}(3) spends both directions, and
\ref{thm:matern-pyramid}(2) the form of a fully covariant family. Not carried: the page's sentence
on the kernels from scale zero, $\widehat\mu_{0,t}(\omega) = e^{-F(\chi(t)|\omega|)}$, which is
the case $s = 0$ of the display; and the page's wider uniqueness, against any real $F'$ on
$[0,\infty)$ with $F'(0) = 0$, which no node uses.
\textbf{Hypotheses} (the node's, each with what its declarations require): $d \ge 1$: same; the family
of probability measures on $\RR^d$: same, \texttt{IsProbabilityMeasure} for every pair
$0 \le s \le t$; $\chi$ a gauge: same; $F \in \mathcal{A}_d$: same (\texttt{classA}); $F \not\equiv 0$:
same, $F(r) \ne 0$ for some $r \ge 0$; the display: same, the transform of $\mu_{s,t}$ at $\xi$ in
terms of $|\xi|$; for the uniqueness, the second pair $(\chi', F')$ with $\chi'$ a gauge and the
display: same, with $F' \in \mathcal{A}_d$ required where the page requires only a real $F'$ on $[0,\infty)$
with $F'(0) = 0$, so stronger.)}
\end{theorem}

The reduction behind it \prog{iso v0.1}{\S~3.2, p.~12} takes the marginal $\pi_*\mu$,
$\pi(x) = \langle e, x\rangle$ for a unit vector $e$. The marginal determines a rotation-invariant
law, $\widehat\mu(\omega) = \widehat{\pi_*\mu}(|\omega|)$, and commutes with convolution and
dilation \prog{iso v0.1}{Lem.~3.4, pp.~12--13}. The marginals of an isotropic family form a family
on the line with the same bijections \prog{iso v0.1}{Lem.~3.5, p.~13}. Every admissible pair is
realized, by the next theorem.

% draft: the sentence "every pair (chi, F) with F in A_d, F not 0, is realised by an isotropic
% family [iso v0.1, Lem. 3.8, p. 14]" (section 2.0), line 81.
\begin{theorem}[every pair is realized]
\label{thm:cited-A34}
\lean{Pyramids.cited_A34}
\leanok
\uses{def:isotropic-cascade-family, def:self-decomposable}
Let $d \ge 1$, let $\chi$ be a gauge and let $F \in \mathcal{A}_d$ with $F \not\equiv 0$. Then
there is an isotropic cascade family $(\mu_{s,t})_{0 \le s \le t}$ in dimension $d$ with
\[
  \widehat\mu_{s,t}(\omega) = \exp\bigl(-[F(\chi(t)|\omega|) - F(\chi(s)|\omega|)]\bigr),
  \qquad 0 \le s \le t,\ \omega \in \RR^d.
\]
% CHANGED (wave 2g, 2026-11-14, as thm:cited-A1 was): regraded [T]. The affine module's formal
% development proves the lemma with no axiom, so the node is a theorem of the formal development
% and not an interface; ledger entry A34 keeps the page anchor and its transcription.
% CHANGED (fidelity review, 2026-11-16, R71): the status line lists the declaration's hypotheses beside the node's (LINKAGE rule 6).
\statusT\quad\emph{(The isotropic paper's realization lemma,\prog{iso v0.1}{Lem.~3.8, p.~14}.
Held as [T]: the affine module's formal development proves it with no axiom
(\texttt{AffineHemigroup.isotropic\_realisation}), the declaration named above is that theorem, and
ledger entry A34 records the page anchor and its transcription. \ref{prop:ladder-knot-kernel}(4) and
\ref{thm:matern-pyramid}(1) read it at $\chi = \mathrm{id}$. The node says nothing about the
uniqueness of the family, which is Theorem~\ref{thm:cited-A31}.
\textbf{Hypotheses} (the node's, each with what its declaration requires): $d \ge 1$: same; $\chi$ a
gauge: same; $F \in \mathcal{A}_d$: same (\texttt{classA}); $F \not\equiv 0$: same, $F(r) \ne 0$ for
some $r \ge 0$.)}
\end{theorem}

% CHANGED (wave 2g, 2026-11-14): the node is proved; the proof of record is the construction of
% the isotropic paper's lemma, as the affine module checks it.
% CHANGED (the theory paper, 2026-11-16): the closing sentence on checking and cited facts shortened; no step changed.
\begin{proof}
\leanok
Let $\mu$ be the self-decomposable law on $\RR^d$ with $\widehat\mu(\omega) = e^{-F(|\omega|)}$,
which $F \in \mathcal{A}_d$ provides. For $0 < c < 1$ it has a residual $\varrho_c$, a probability
measure with $\widehat\mu(\omega) = \widehat\mu(c\omega)\,\widehat\varrho_c(\omega)$; put
$\varrho_0 := \mu$ and $\varrho_1 := \delta_0$. Since $\widehat\mu$ has no zero and $F(0) = 0$,
\[
  \widehat\varrho_c(\omega) = \exp\bigl(-[F(|\omega|) - F(c|\omega|)]\bigr), \qquad 0 \le c \le 1 .
\]
A gauge has $\chi(0) = 0$ and $\chi(t) > 0$ for $t > 0$. For $0 \le s \le t$ with $t > 0$ let
$\mu_{s,t}$ be the image of $\varrho_{\chi(s)/\chi(t)}$ under $x \mapsto \chi(t)x$, and let
$\mu_{0,0} := \delta_0$. The transform of $\mu_{s,t}$ at $\omega$ is that of $\varrho_{\chi(s)/\chi(t)}$
at $\chi(t)\omega$, which is the display.

The display gives the axioms of Definition~\ref{def:isotropic-cascade-family}. (I1): the exponents
of $\mu_{s,t}$ and $\mu_{r,s}$ add up to that of $\mu_{r,t}$, and a law is determined by its
transform. (I2): $F$ and $\chi$ are continuous. (I3): the transform depends on $\omega$ only through
$|\omega|$. (I4): with $S_\lambda := \chi^{-1}(\lambda\chi(\cdot))$, an increasing bijection of
$[0,\infty)$, the transform of $(\lambda I)_*\mu_{s,t}$ at $\omega$ is the display at
$(S_\lambda s, S_\lambda t)$. (ND): if $\mu_{s,t} = \delta_0$ with $s < t$, then $F(cr) = F(r)$ for
every $r \ge 0$, with $c := \chi(s)/\chi(t) < 1$; iterating, $F(r) = F(c^nr) \to F(0) = 0$, against
$F \not\equiv 0$.

Every step is machine-checked, in the affine paper's formal development, with no cited fact.
\end{proof}

Chapter~3 keeps (I1), (I2) and (ND), treats the rotation invariance (I3) as an optional property,
and requires (I4) at one dilation (Definition~\ref{def:family-dimension-d}). The class $\mathcal{A}_d$
has a polar description, by the \Levy{} measures of its laws.

% draft: the row "[iso v0.1, Def. 3.10] p. 15, the isotropic origin data and the polar profile" of
% the table in section 2.0 (line 109); the definition as the draft reads it at lines 418 and 553.
% CHANGED: the two integrability conditions of the definition, which the draft's prose at line 418
% omits, are stated (the definition has them and the draft uses them at line 553;
% VERBATIM-selection-iso-affine.md, Def. 3.10). The page's last two sentences (the profile at the
% origin and subordinated data) are not restated: no node of this blueprint uses them.
% CHANGED (fidelity review, 2026-11-15, R3): the tag gains OriginDatum.exponentL, the
% extended-real exponent of which OriginDatum.exponent is the real part, and
% OriginDatum.exponentL_ne_top, the lemma (for d >= 1) that makes that real part the exponent.
% CHANGED (the theory paper's blueprint round, 2026-11-16, review of Section 2, NOTATION): the
% Gaussian coefficient a of the isotropic paper, which enters as a|xi|^2/2, is written a_0, as
% chapter 3 writes the origin datum, apart from the a of (2.0), which enters as a omega^2. A rename.
\begin{definition}[isotropic origin data]
\label{def:isotropic-origin-data}
\lean{AffineHemigroup.OriginDatum, AffineHemigroup.OriginDatum.levyMeasure,
AffineHemigroup.OriginDatum.exponentL, AffineHemigroup.OriginDatum.exponentL_ne_top,
AffineHemigroup.OriginDatum.exponent, AffineHemigroup.IsLawOf, AffineHemigroup.sphereArea}\leanok
\uses{def:self-decomposable}
Let $d \ge 1$, write $\sigma_{d-1}$ for the surface measure of the unit sphere
$S^{d-1} \subset \RR^d$ and $|S^{d-1}| := \sigma_{d-1}(S^{d-1})$. An \emph{isotropic origin datum
in dimension $d$} is a pair $(a_0, k_d)$ with $a_0 \ge 0$ and $k_d : (0,\infty) \to [0,\infty)$
nonincreasing, subject to
\[
  \int_0^1 k_d(r)\,r\,dr < \infty \qquad\text{and}\qquad \int_1^\infty k_d(r)\,\frac{dr}{r} < \infty .
\]
Its \emph{\Levy{} measure} is the rotation-invariant measure $\nu_d$ on $\RR^d \setminus \{0\}$
with radial density $f_d(r) = k_d(r)\bigl(|S^{d-1}|\,r^d\bigr)^{-1}$, so that
\[
  \int_{\RR^d} g\,d\nu_d = \frac{1}{|S^{d-1}|}\int_{S^{d-1}}\sigma_{d-1}(d\eta)
  \int_0^\infty g(r\eta)\,k_d(r)\,\frac{dr}{r}
\]
for every nonnegative Borel $g$ vanishing at the origin. Its \emph{exponent} is
\[
  F_d(\xi) = \tfrac12 a_0|\xi|^2 + \int_{\RR^d}\bigl(1 - \cos\langle\xi, x\rangle\bigr)\,\nu_d(dx),
  \qquad \xi \in \RR^d,
\]
and its \emph{law} is the symmetric infinitely divisible law $\mu$ on $\RR^d$ with
$\widehat\mu = e^{-F_d}$. The function $k_d$ is the datum's \emph{polar profile}.
\statusT\quad\emph{(A definition, \prog{iso v0.1}{Def.~3.10, p.~15}, displays (3.2) and (3.3), whose
$a$ is written $a_0$ here. Jumps of length between $r_1$ and $r_2$ occur at the rate
$\int_{r_1}^{r_2}k_d(r)\,dr/r$, as on the line. The Gaussian coefficient $a_0$ enters as
$\tfrac12 a_0$, in Sato's normalization ($A^{\mathrm S} = a_0I$), against $a$ on the line,
(\ref{eq:admissible-exponent}); Theorem~\ref{thm:cited-A35}(2) relates the two. That the law exists is
Theorem~\ref{thm:cited-A39}(iii), the two conditions being the integrability (8.2) of $\nu_d$. It asserts
nothing further and cites no ledger entry.)}
\end{definition}

% draft: the row "[iso v0.1, Prop. 3.11] p. 15, the polar form of A_d" of the table in section 2.0
% (line 110), as the draft reads it at lines 418, 428, 500 and 553.
% CHANGED: the page is pp. 15-16, clause (1) on p. 15 and clause (2) on p. 16
% (VERBATIM-selection-iso-affine.md, A35). Clause (3) of the proposition, the classification in
% polar form (plan id A36), is not an entry: no node spends it.
% CHANGED (the theory paper's blueprint round, 2026-11-16, review of Section 2, NOTATION): the
% origin data are written (a_0, k_d) and (a_0', k_1), as in def:isotropic-origin-data, and (2) names
% the coefficient a of (2.0). A rename; the content is the page's.
\begin{theorem}[the polar form of $\mathcal{A}_d$]
\label{thm:cited-A35}
\lean{Pyramids.cited_A35_reading, Pyramids.cited_A35_unique, Pyramids.cited_A35_admissible,
Pyramids.cited_A35_classA}
\leanok
\uses{def:isotropic-origin-data, def:self-decomposable, thm:cited-A39, thm:cited-A43,
thm:cited-A45}
Let $d \ge 1$ and $F : [0,\infty) \to [0,\infty)$.
\begin{enumerate}
\item $\xi \mapsto F(|\xi|)$ is the exponent of a self-decomposable law on $\RR^d$ if and only if
it is the exponent $F_d$ of an isotropic origin datum $(a_0, k_d)$ in dimension $d$
(Definition~\ref{def:isotropic-origin-data}):
\[
  F(|\xi|) = \tfrac12 a_0|\xi|^2 + \int_{\RR^d}\bigl(1 - \cos\langle\xi, x\rangle\bigr)\,
  \frac{k_d(|x|)}{|S^{d-1}|\,|x|^d}\,dx, \qquad \xi \in \RR^d.
\]
The datum is then unique, $k_d$ up to modification on a Lebesgue-null set.
\item If this holds, $F$ is admissible, of the form (\ref{eq:admissible-exponent}) with Gaussian
coefficient $a = \tfrac12 a_0'$ and profile $k = k_1$ for its datum $(a_0', k_1)$ in dimension $1$. So
$\mathcal{A}_d$ is the set of $F$ for which $F(|\cdot|)$ is the exponent of an isotropic origin
datum in dimension $d$.
\end{enumerate}
% CHANGED (fidelity review, 2026-11-16, R71): the status line lists the declarations' hypotheses beside the node's (LINKAGE rule 6).
\statusA\quad\ledger{A35}\quad\emph{(Cited from \prog{iso v0.1}{Prop.~3.11(1)--(2), pp.~15--16}:
clause (1) on p.~15, clause (2) on p.~16. Admitted as a Lean axiom by the author on 2026-11-14,
and found to need none of its own: the affine module's formal development proves both clauses
(\texttt{AffineHemigroup.polar\_reading}, \texttt{polar\_reading\_unique},
\texttt{polar\_reading\_admissible}, \texttt{classA\_iff\_originDatum}) from
Theorem~\ref{thm:cited-A39}, the first clause of Theorem~\ref{thm:cited-A43} and Theorem~\ref{thm:cited-A45}, and the
declarations named above are those theorems. \textbf{Assignment.} \ledger{A35} carries both clauses. \ref{thm:matern-pyramid}(1) and (5)
spend the membership in $\mathcal{A}_d$ of the exponent of an isotropic origin datum, which needs
both clauses since $\mathcal{A}_d$ is defined among the admissible exponents;
\ref{prop:ladder-knot-kernel}(4) spends clauses (1) and (2), in (4)(b): the equivalence of (1)
with the uniqueness of the datum, and the admissibility of (2).
Not carried: clause (3) of the proposition, the classification in polar form, which no node spends;
the strict filtration $\mathcal{A}_{d+1} \subsetneq \mathcal{A}_d$ of the isotropic paper's
Prop.~3.22, which chapter~3 cites in prose.
\textbf{Hypotheses} (the node's, each with what its declarations require): $d \ge 1$: same, in all four;
$F : [0,\infty) \to [0,\infty)$: same, $F(r) \ge 0$ for every $r \ge 0$ (the declarations take $F$ on
$\RR$ and constrain it on $[0,\infty)$ only); for the uniqueness of the datum, two origin data with
$F(|\xi|)$ their exponent: same; for clause (2), that $F(|\cdot|)$ is the exponent of a
self-decomposable law: same, a probability measure that is self-decomposable with transform
$e^{-F(|\xi|)}$.)}
\end{theorem}

\subsection{The generator form}
\label{subsec:cited-generator}

Let $F$ be admissible with data $(a, k)$, put $\kappa^{\mathrm C} := -dk$, a measure on
$(0,\infty)$, and let $u(t,\cdot) = \mu_{0,t} * f$ be the scale space of a Schwartz function
$f \in \mathcal{S}(\RR)$ in the canonical gauge. Then $\partial_t u = \mathcal{H}_t u$ with
\begin{equation}
\label{eq:generator-form}
\tag{2.0$''$}
  \mathcal{H}_t g(x) = 2at\,g''(x) + \frac1t\int_{(0,\infty)}\delta^2_{tv}g(x)\,
  \kappa^{\mathrm C}(dv), \qquad
  \delta^2_h g(x) := \tfrac12\bigl(g(x - h) + g(x + h)\bigr) - g(x)
\end{equation}
\prog{V-cone v0.1}{Prop.~6.2, display (6.2), p.~56}. The Gaussian member has
$\mathcal{H}_t = 2at\,\partial_x^2$ \prog{V-cone v0.1}{Prop.~6.4(1), p.~58}. Write
$\Cin(z) := \int_0^z(1 - \cos v)\,dv/v$, the entire cosine integral. The \emph{Cin ray of step
$\tau$} is the exponent $\Cin(\tau\omega)$, which is (\ref{eq:admissible-exponent}) with $a = 0$ and
the profile $k = \mathbf{1}_{(0,\tau)}$ \prog{V-cone v0.1}{Lemma~3.1(1), p.~13}. Its generator is a
single symmetric second difference.

% draft: the sentence "Its generator is a single symmetric second difference, H_t = t^{-1}
% delta^2_{t tau} [V-cone v0.1, Prop. 6.4(5), p. 58]" (section 2.0), line 85.
% CHANGED: the statement reads the clause through Prop. 6.2 of the same module, where the generator
% and its domain are defined; the clause itself prints no domain (VERBATIM-V-cone.md, A13).
\begin{theorem}[the generator of a Cin ray]
\label{thm:cited-A13}
\lean{Pyramids.cited_A13}
\notready
Let $\tau > 0$, and let $F$ be the admissible exponent with data $(0, \mathbf{1}_{(0,\tau)})$, the
Cin ray of step $\tau$. In the canonical gauge its generator is a single symmetric second
difference: for $f \in \mathcal{S}(\RR)$ the scale space $u(t,\cdot) = \mu_{0,t} * f$ satisfies,
for every $t > 0$,
\[
  \partial_t u(t,\cdot) = \frac1t\,\delta^2_{t\tau}u(t,\cdot), \qquad
  \delta^2_h g(x) = \tfrac12\bigl(g(x - h) + g(x + h)\bigr) - g(x).
\]
\statusA\quad\ledger{A13}\quad\emph{(Cited from \prog{V-cone v0.1}{Prop.~6.4(5), p.~58}, ``Cin
ray of step $\tau$: $\mathcal{A}_t = \frac1t\delta^2_{t\tau}$'', with the generator $\mathcal{A}_t$
and the signal class $\mathcal{S}(\RR)$ of the cone module's Prop.~6.2 and Def.~6.1, p.~56. Its statement is typed in the skeleton
(\texttt{PyramidsSkeleton.Ch02CitedFacts}) and is not an axiom; admitting it is the author's decision at phase~5. \textbf{Assignment.} \ledger{A13}
carries the whole statement. No node spends it as a step of a proof: the theory on the line refers
to it (\S\ref{subsec:line-strata}) to say that a grid's second difference is the generator of a Cin
ray held at a fixed step. Not carried: the generator form (\ref{eq:generator-form}) for other
admissible exponents, the cone module's Prop.~6.2, which is cited in prose above; and the value
$F(\omega) = \Cin(\tau\omega)$ of the ray, the cone module's Lemma~3.1(1).)}
\end{theorem}

So (\ref{eq:generator-form}) writes every admissible evolution as a positive superposition of these
\emph{atomic generators}, with the weights $(a, \kappa^{\mathrm C})$. The same weights give $F$
itself, and the Gaussian ray and the Cin rays are the extreme rays of the cone of admissible
exponents. The cone module proves it, and it is stated here for the reader.

% CHANGED (2026-11-23, review round 3 batch C, item 15, decision D3 of
% records/pyramids/PLAN-review-response-theory-round3.md): new interface node, ledger entry A15
% (the plan id the blueprint brief gave V-cone Prop. 3.3, unspent until now). The cone module's
% Prop. 3.3 was cited in the prose above ("Every admissible exponent is a unique superposition ...
% its Choquet representation [V-cone v0.1, Prop. 3.3, pp. 15-16]"); it is restated at the page's
% letter (records/pyramids/VERBATIM-V-cone.md, A15; reviews/LIBRARIAN-r3.md, "For batch C"), as
% the six other unread facts are. Its statement is typed in the skeleton.
\begin{theorem}[the Choquet structure of the admissible cone]
\label{thm:cited-A15}
\lean{Pyramids.cited_A15_admissible, Pyramids.cited_A15_injective, Pyramids.cited_A15_inverse,
Pyramids.cited_A15_extreme}
\notready
The superposition map $(a, \varpi) \mapsto F$,
\[
  F(\omega) = a\omega^2 + \int_{(0,\infty)}\Cin(\tau\omega)\,\varpi(d\tau), \qquad \omega \in \RR,
\]
is a bijection from the pairs $(a, \varpi)$, $a \ge 0$ and $\varpi$ a positive Borel measure on
$(0,\infty)$ with $\int\min\bigl(\tau^2, 1 + \log_+\tau\bigr)\,\varpi(d\tau) < \infty$, onto the
admissible exponents. Its inverse takes the exponent with data $(a, k)$ to $(a, \varpi)$ with
$\varpi = -dk$, that is $k(x) = \varpi((x,\infty))$ for almost every $x > 0$. Consequently the
extreme rays of the cone of admissible exponents are the Gaussian ray $\omega \mapsto \omega^2$
and the Cin rays $\omega \mapsto \Cin(\tau\omega)$, $\tau > 0$, and every admissible exponent is a
unique superposition of Cin rays and the Gaussian term $a\omega^2$.
\statusA\quad\ledger{A15}\quad\emph{(Cited from \prog{V-cone v0.1}{Prop.~3.3, pp.~15--16}: the
statement opens at the foot of p.~15 and ends on p.~16, the extreme rays inside it. Stated for the
reader. Its statement is typed in the skeleton (\texttt{PyramidsSkeleton.Ch02CitedFacts}) and is
not an axiom; admitting it is the author's decision at phase~5. \textbf{Assignment.} \ledger{A15}
carries the whole statement. No node spends it as a step of a proof: the generator form above and
the cone of the continuous stratum after Proposition~\ref{prop:strata-cylinder-coordinates} refer
to it. Not carried: the page's word ``linear'', the additivity and positive homogeneity of the map,
which the formula shows and no node reads; the folded \Levy{} density
$x^{-1}\mathbf{1}_{(0,\tau)}(x)$ of the Cin ray of step $\tau$, which the page names and which is
the cone module's Lemma~3.1(1), cited in the prose above. The page's $k(x) = \varpi((x,\infty))$
is read for almost every $x$: the exponent fixes the profile $k$ only up to a null set.
\textbf{Hypotheses} (the node's, each with what its declarations require): $a \ge 0$ and $\varpi$ a
measure on $(0,\infty)$ with the integrability: same (\texttt{IsChoquetPair}), $\varpi$ typed as a
measure on $\RR$ with no mass on $(-\infty, 0]$ and $\log_+$ as \texttt{Real.posLog}; the cone: same,
the functions with \texttt{IsAdmissibleExponent}; the data $(a, k)$: same, an \texttt{SDProfile},
whose field $k(0) = 0$ is a normalization; the extreme rays: same, \texttt{IsExtremeRayIn} on the
cone with nonnegative scalars.)}
\end{theorem}

The theory on the line also uses the \Matern{} and Thorin clauses of the generator proposition, which sum
(\ref{eq:generator-form}) over Laplace kernels \prog{V-cone v0.1}{Prop.~6.4(3)--(4), p.~58}; the
Thorin clause holds for $a = 0$, at every $g$ that is twice continuously differentiable with $g$ and
$g''$ bounded.

\subsection{Other statements of the released modules}
\label{subsec:cited-released}

The later chapters use the following statements of the released modules. Each is stated where it
is used, or below when a proof spends it.
\begin{itemize}
\item \prog{V-line v0.1}{Lemma~6.1, p.~29}: covariance in Fourier form.
\item \prog{V-line v0.1}{Remark~6.7, pp.~32--33}: the gauge freedom of the scale coordinate.
\item \prog{V-cone v0.1}{\S~4, p.~24}: the subordination bridge.
\item \prog{V-cone v0.1}{Prop.~5.7, p.~48}: the symmetric Thorin subclass.
\item \prog{V-cone v0.1}{Prop.~7.2, pp.~60--61}: the \Matern{} cascade by forward--backward
first-order sections (Theorem~\ref{thm:cited-A19}).
\item \prog{V-cone v0.1}{Prop.~7.5(2), p.~65}: rational increments and the pyramid they
interpolate (Theorem~\ref{thm:cited-A20}).
\item \prog{V-cone v0.1}{Remark~7.4, p.~63}: which recursive Gaussians are scale spaces.
\item \prog{V-cone v0.1}{Prop.~7.8, p.~69}: the closure of the families realized by lag sections.
\item \prog{V-cone v0.1}{Remarks~7.10 and 7.11, p.~72}: the local equation of the Student-$t$
corner cannot be marched; the costs on the line, per sample and per knot.
\item \prog{V-selection v0.1}{Thm.~4.2, pp.~30--31}: non-enhancement at global extrema for every family,
at local extrema on the Gaussian ray only; and its Cor.~4.4, p.~33: off the Gaussian ray the scale
space enhances a local maximum.
\item \prog{V-selection v0.1}{Prop.~4.5, p.~35}: non-creation from the signal; the \Polya{}
frequency members (Theorems~\ref{thm:cited-A22} and~\ref{thm:cited-A23}).
\item \prog{V-selection v0.1}{Thm.~4.10, p.~43}: scale-monotone non-creation characterizes the
Gaussian.
\item \prog{V-selection v0.1}{Def.~5.1, p.~54}: the spatial jet on the line.
\item \prog{V-selection v0.1}{Props.~5.6 and 5.7, pp.~60--62}: the \Matern{} jet by subtraction,
the Helmholtz identity, and the cascade outputs as the jet.
\item \prog{V-selection v0.1}{Example~5.9 and App.~B, pp.~63--65 and 73--76}: the jet from the
taps, on one signal.
\item \prog{V-selection v0.1}{\S~6, p.~69}: diffusion creates extrema in the plane.
\item \prog{iso v0.1}{Prop.~3.12, pp.~16--17}: the subordination bridge in $\RR^d$; the
\Matern{}, stable and Student-$t$ kernels.
\item \prog{iso v0.1}{Prop.~3.22, p.~27}: the polar profile one dimension down;
$\mathcal{A}_{d+1} \subsetneq \mathcal{A}_d$ for every $d$.
\item \prog{affine v0.1}{Prop.~4.4, p.~38}: the affine orbits of isotropic laws.
\item \prog{affine v0.1}{Prop.~F.1, p.~93}: the frame families.
\end{itemize}

% draft: the row "[V-cone v0.1, Prop. 7.2] p. 60" of the table in section 2.0 (line 95).
% CHANGED: stated at the kernels of the sections, which do not depend on the sign convention of
% the transform; the page's complex factorization of (7.1) follows the cone module's convention and
% is not restated. The cost sentence of the page is not restated (no node uses it).
% CHANGED (the theory paper's blueprint round, 2026-11-16, review of Section 2, NOTATION): the kernels
% Exp_+-(t) and Lap_t are written e^+-_t and L_t, the names of prop:recursion-helmholtz(1), so that
% the continuum Laplace law has one name. A rename.
\begin{theorem}[the \Matern{} cascade by first-order sections]
\label{thm:cited-A19}
\lean{Pyramids.cited_A19_transform, Pyramids.cited_A19_sections, Pyramids.cited_A19_forward,
Pyramids.cited_A19_backward, Pyramids.cited_A19_gamma_one}
\notready
Let $\gamma \ge 1$ be an integer, and consider the \Matern{} family on the line in the canonical
gauge, $e^{-F(t\omega)} = (1 + t^2\omega^2)^{-\gamma}$. For $t > 0$ let $e^\pm_t$ be
the one-sided exponential kernels $t^{-1}e^{-|x|/t}$ carried by $\pm(0,\infty)$, and
$L_t$ the Laplace kernel $(2t)^{-1}e^{-|x|/t}$. Let $0 < t_k < t_{k+1}$ be adjacent knots
and $r := t_k/t_{k+1} \in (0,1)$. Then the increment has transform
\[
  \widehat\mu_{t_k, t_{k+1}}(\omega) = \Bigl(\frac{1 + t_k^2\omega^2}{1 + t_{k+1}^2\omega^2}\Bigr)^{\gamma},
\]
and it is a cascade of $\gamma$ identical causal first-order sections and $\gamma$ identical
anticausal ones. Each causal section is the convolution with
$r\,\delta_0 + (1 - r)\,e^+_{t_{k+1}}$, a mixture of no displacement and a one-sided
exponential displacement, with the one-state realization
\[
  w' = \frac{y - w}{t_{k+1}}, \qquad v = r\,y + (1 - r)\,w,
\]
run forward in $x$ (input $y$, state $w$, output $v$, the prime the derivative in $x$), the
solution $w$ being the one in $L^1$, with $w(-\infty) = 0$. Each anticausal
section is its mirror image, run backward. For $\gamma = 1$ the two-sided increment is
\[
  r^2\delta_0 + r(1 - r)\bigl(e^+_{t_{k+1}} + e^-_{t_{k+1}}\bigr)
  + (1 - r)^2\,L_{t_{k+1}} .
\]
\statusA\quad\ledger{A19}\quad\emph{(Cited from \prog{V-cone v0.1}{Prop.~7.2, pp.~60--61}: the
statement opens at the foot of p.~60 with (7.1), and the realization and (7.2) are on p.~61. Its statement is typed in the skeleton
(\texttt{PyramidsSkeleton.Ch02CitedFacts}) and is not an axiom; admitting it is the author's decision at phase~5. \textbf{Assignment.} \ledger{A19}
carries the whole statement. No node spends it as a step of a proof: the theory on the line refers
to it (\S\ref{subsec:sampled-laplace}) for how the \Matern{} family is computed on the line.
Not carried: the rung from scale zero, $t_k = 0$, whose increment is
$L_{t_{k+1}} = e^+_{t_{k+1}} * e^-_{t_{k+1}}$; the page states the
convex-combination clause for $0 < t_k/t_{k+1} < 1$ only, and the identity
$e^+_t * e^-_t = L_t$ is in its proof, p.~61. Nor the cost of a
sampled cascade, which the page states and no node uses.)}
\end{theorem}

% draft: the quotation of Remark 2.22 (section 2.6), line 293, and the row of the table in
% section 2.0. The draft's quotation is a paraphrase; this node is at the page's words
% (VERBATIM-V-cone.md, A20).
% CHANGED: the page is p. 65 (the draft's "pp. 60-72" is wrong; VERBATIM-V-cone.md, A20).
\begin{theorem}[rational increments and the pyramid they interpolate]
\label{thm:cited-A20}
\lean{Pyramids.cited_A20_unique, Pyramids.cited_A20_product, Pyramids.cited_A20_profile,
Pyramids.cited_A20_admissible, Pyramids.cited_A20_admissible_of_antitone}
\notready
Let $R$ be a symmetric rational function that is the transform of an infinitely divisible law on
$\RR$, with folded \Levy{} profile $k_R$. Then for each $q > 1$ there is exactly one function $F$
that is continuous with $F(0) = 0$ and satisfies $F(q\omega) - F(\omega) = -\log R(\omega)$,
namely
\[
  e^{-F(\omega)} = \prod_{m \ge 1}R(q^{-m}\omega), \qquad\text{with folded profile}\qquad
  k(x) = \sum_{m \ge 1}k_R(q^m x);
\]
% CHANGED (skeleton review, m2): "is nonincreasing" reads "has a nonincreasing version", the
% reading of the page's own proof, "the monotonicity is monotonicity of a representative"
% (VERBATIM-V-cone.md, A20); read pointwise the sentence fails for a profile changed on a null set.
it is admissible exactly when this $k$ has a nonincreasing version ($k_R$, and with it $k$, being
fixed only up to a null set), in particular whenever $k_R$ is nonincreasing. The
\Matern{} pyramid is the case $R(\omega) = (1 + \omega^2)^{-1}$.
\statusA\quad\ledger{A20}\quad\emph{(Cited from \prog{V-cone v0.1}{Prop.~7.5(2), p.~65}, statement
and proof on that page. A folded profile $k_R$ is a function with folded \Levy{} measure
$k_R(x)\,dx/x$ on $(0,\infty)$. The function $F$ has Gaussian coefficient $0$: the proof on the same
page shows that $-\log R$ has none and that $F$ is of the form (2.3) ``at the pair $(0,k)$''. So
``with folded profile $k$'' is read as
$F(\omega) = \int_0^\infty(1 - \cos\omega x)\,k(x)\,x^{-1}\,dx$, the form the typed statement
asserts. Its statement is typed in the
skeleton (\texttt{PyramidsSkeleton.Ch02CitedFacts}) and is not an axiom; admitting it is the
author's decision at phase~5.
\textbf{Assignment.} \ledger{A20} carries the whole statement. No node spends it as a step of a
proof: the theory on the line compares it with the lattice statement
\ref{prop:grid-pyramids-infinitely-divisible} (\S\ref{subsec:grid-pyramids-line}, the lead and the paragraph
after that proposition), which proves its own. Not carried: the proof's remark that for real poles $k_R$
is a finite combination of exponentials, which the statement does not print; and clause (1) of the
proposition.)}
\end{theorem}

% draft: the row "[V-selection v0.1, Prop. 4.5] p. 35" of the table in section 2.0 (line 102), as
% the theory on the line uses it (draft lines 229, 235, 243, 247).
% CHANGED: narrowed to the "if" of the page's equivalence, the direction the nodes spend
% (VERBATIM-selection-iso-affine.md, A22). The hypothesis "admissible" is kept.
\begin{theorem}[the \Polya{} frequency members]
\label{thm:cited-A22}
\lean{Pyramids.cited_A22_variation, Pyramids.cited_A22_polya}
\leanok
For a function $g$ on $\RR$ write $S^-(g)$ for the number of sign changes of $g$, in the essential
sense.
% CHANGED (skeleton review, m1): F is not identically zero, the presupposition of "an admissible
% family" on the page (VERBATIM-selection-iso-affine.md, A22); the kernel of F = 0 is the point
% mass, which is no Polya frequency density. Every spender's exponent is not identically zero.
Let $F \not\equiv 0$ be an admissible exponent of the form
\[
  F(\omega) = a\,\omega^2 + \sum_i\log\Bigl(1 + \frac{\omega^2}{\theta_i^2}\Bigr), \qquad
  a \ge 0,\quad \theta_i > 0,\quad \sum_i\theta_i^{-2} < \infty,
\]
over countably many $\theta_i$. Then the similarity family of $F$, with kernels
$\widehat\mu_{0,t} = e^{-F(t\,\cdot)}$, has \Polya{} frequency kernels, so that convolution with
$\mu_{0,t}$ is variation-diminishing: $S^-(\mu_{0,t} * g) \le S^-(g)$ for every $t > 0$ and every
integrable or bounded $g$.
% CHANGED (fidelity review, 2026-11-15, R26/R50/R69): the status line lists the declarations' hypotheses beside the node's (LINKAGE rule 6).
\statusA\quad\ledger{A22}\quad\emph{(Cited from \prog{V-selection v0.1}{Prop.~4.5(1), p.~35}, the
``if'' of its equivalence. Admitted as two Lean axioms by the author on 2026-11-15, one per
clause, both for an admissible $F \not\equiv 0$: the \Polya{} frequency kernels and the variation
diminution. \textbf{Assignment.} \ledger{A22} carries the statement. \ref{thm:matern-pyramid-line-noncreation}
spends it for $F_q$ and for $\log(1 + \omega^2)$, and \ref{prop:ladder-non-creation-line}(1) for
the exponent $G$ of its proof, whose form and admissibility that node proves. Not carried: the
``only if'' of the page, that no other admissible family has \Polya{} frequency kernels; and the
page's list of members and non-members.
\textbf{Hypotheses} (the node's, each with what its declarations require): $F$ admissible: same
(\texttt{IsAdmissibleExponent}, a profile of the form (\ref{eq:admissible-exponent})); $F \not\equiv 0$:
same; $a \ge 0$, $\theta_i > 0$, $\sum_i\theta_i^{-2} < \infty$, countably many $\theta_i$ and the
displayed form: same (a countable index type, the sum a \texttt{tsum} with its summability required);
the kernels of the similarity family: same, any probability measures $\mu_t$ on $\RR$ with
$\widehat\mu_t = e^{-F(t\,\cdot)}$ for $t > 0$, the transform determining the law; $t > 0$: same;
$g$ integrable: same; $g$ bounded: stronger, \texttt{cited\_A22\_variation} also requires it
(a.e.-strongly) measurable, which the convolution presupposes and the node leaves implicit.)}
\end{theorem}

% draft: the same row (line 102), the counting clause, as used at lines 243 and 247.
\begin{theorem}[the counting clause]
\label{thm:cited-A23}
\lean{Pyramids.cited_A23}
% CHANGED (wave 3b, 2026-11-15): proved from A22's variation-diminishing clause by the page's own
% derivation (p. 39); statement and the declaration's type unchanged. The entry is not admitted
% and needs no axiom of its own (ADR-0005).
% CHANGED (fifth blind pass, 2026-11-15, flags 1, 3): the derivation moves out of the annotation into
% a proof environment below; the Assignment says that A23 grounds the node and that the proof
% spends the variation-diminishing clause of A22. It read "A23 carries the whole statement",
% beside "not admitted ... proved from thm:cited-A22".
\leanok
\uses{thm:cited-A22}
For a family as in Theorem~\ref{thm:cited-A22} and $f$ the primitive of an integrable $f'$, the
scale space $u(t,\cdot) = \mu_{0,t} * f$ has, at every scale $t > 0$, no more crossings of any
value $c$ than the signal and no more changes of monotonicity:
% CHANGED (register pass, 2026-11-16): "zero crossings than the signal at every level" read as a
% pyramid level; the wording now names the value c of the display. The claim is unchanged.
\[
  S^-\bigl(u(t,\cdot) - c\bigr) \le S^-(f - c) \quad (c \in \RR), \qquad
  S^-\bigl(\partial_x u(t,\cdot)\bigr) \le S^-(f').
\]
% CHANGED (fidelity review, 2026-11-16, R71): the status line lists the declaration's hypotheses beside the node's (LINKAGE rule 6), and says that the declaration does not require F not identically zero.
\statusA\quad\ledger{A23}\quad\emph{(Cited from \prog{V-selection v0.1}{Prop.~4.5(2), p.~35}; its
derivation from clause (1), applied to the bounded $f - c$ and to $f'$ with
$\partial_x u(t,\cdot) = \mu_{0,t} * f'$, is on p.~39. This citation is the node's grounding. Not
admitted, and found to need no axiom of its own: the proof below derives the statement from
Theorem~\ref{thm:cited-A22} by that derivation. \textbf{Assignment.} \ledger{A23} is where the statement
is printed, and grounds the node; what the proof spends from the literature is the
variation-diminishing clause of \ledger{A22} (Theorem~\ref{thm:cited-A22}), applied to $f - c$ and to $f'$.
\ref{prop:ladder-non-creation-line}(2) spends the statement for the similarity family of the exponent $G$ of
its proof, applied to the signal at the rung. \ref{thm:matern-pyramid}(3) spends it for the
line's \Matern{} pyramid, $a = 0$ and $\theta_i = q^i/c$, applied to a version of the projected
signal, for the count from scale zero. Not carried: any statement about a signal that is not
the primitive of an integrable function.
\textbf{Hypotheses} (the node's, each with what its declaration requires): $F$ admissible, of the form
of Theorem~\ref{thm:cited-A22} with $a \ge 0$, $\theta_i > 0$, $\sum_i\theta_i^{-2} < \infty$ and countably
many $\theta_i$: same (as for Theorem~\ref{thm:cited-A22}); $F \not\equiv 0$, which ``a family as in
Theorem~\ref{thm:cited-A22}'' carries: weaker, \texttt{cited\_A23} does not require it, the case
$F \equiv 0$, where the kernel is the point mass, being covered by a separate lemma, so the
declaration is the stronger statement; the kernels of the family: same, probability measures with
$\widehat\mu_t = e^{-F(t\,\cdot)}$ for $t > 0$; $f'$ integrable: same; $f$ its primitive: same,
$f = C + \int_0^x f'$ for some constant $C$; $t > 0$: same.)}
\end{theorem}

% CHANGED (fifth blind pass, 2026-11-15, flags 2, 3): new proof environment, the derivation formerly
% printed inside the annotation, as thm:cited-A9, thm:cited-A4 and thm:cited-A44 have since the
% fourth pass. The sentence on F = 0, a case the statement excludes, is dropped.
\begin{proof}
\leanok
\uses{thm:cited-A22}
The primitive $f = C + \int_0^x f'$ is bounded by $|C| + \|f'\|_1$ and continuous, so
$f - c$ is a bounded measurable test function and $\mu_{0,t} * f - c = \mu_{0,t} * (f - c)$ at
every point. By Fubini's theorem on $[0, x] \times \mu_{0,t}$,
$u(t, x) = u(t, 0) + \int_0^x (\mu_{0,t} * f')(z)\,dz$ at every $x$, with $\mu_{0,t} * f'$ integrable;
by the Lebesgue differentiation theorem $\partial_x u(t,\cdot) = \mu_{0,t} * f'$ almost everywhere,
and $S^-$, a count in the essential sense, does not see the null set. The variation-diminishing
clause of Theorem~\ref{thm:cited-A22}, applied to the bounded $f - c$ and to the integrable $f'$,
gives the two inequalities.
\end{proof}

\subsection{Notation}
\label{subsec:cited-notation}

This paper writes the law with transform $(1 + r^2|\omega|^2)^{-\gamma}$ as the \Matern{} law
\emph{of index $\gamma$ at range $r$}. The isotropic paper's \Matern{} kernel of index $\gamma$ at
scale $t$ has transform $(1 + t^2|\omega|^2/2)^{-\gamma}$ \prog{iso v0.1}{Prop.~3.12(2), p.~16},
so it is the law at range $t/\sqrt2$. The cone module's \Matern{} kernel at scale $t$ has transform
$(1 + t^2\omega^2)^{-\gamma}$ \prog{V-cone v0.1}{Prop.~5.2(2), p.~43}, the law at range $t$.
% CHANGED: the draft cites the cone module's Prop. 6.4(3), p. 58, for the transform of its Matern
% kernel; that clause is the generator, and the transform is Prop. 5.2(2), p. 43 (VERBATIM-V-cone.md,
% A12 and A17).

The transform of a measure is $\widehat\mu(\omega) = \int e^{i\langle\omega, x\rangle}\mu(dx)$. For
a linear map $A$ of $\RR^d$, $A_*\mu$ is the image of $\mu$ under $x \mapsto Ax$, and
$D_c := (cI)_*$ for $c > 0$, so that $\widehat{D_c\mu}(\omega) = \widehat\mu(c\omega)$. On the line
the dilation of measures is written $T_c$: for a measure $\mu$ on $\RR$ and $c > 0$, $T_c\mu$ is the
image of $\mu$ under $x \mapsto cx$, $(T_c\mu)(B) = \mu(c^{-1}B)$ for Borel sets $B \subseteq \RR$,
so $T_c = D_c$ at $d = 1$, as in [\sscite{sato1999levy}]. The same notation is used for a measure on
$(0,\infty)$, such as a folded \Levy{} measure, since $x \mapsto cx$ maps $(0,\infty)$ onto itself.
The dilation of functions $\mathcal{D}_\lambda$ is the one of
(A8), mass-preserving. The ratio of a pyramid is $q > 1$, and $\ell := \log q$. The \Levy{} data of
a law on $\RR^d$ are a Gaussian part $A$, entering the exponent as $\langle\omega, A\omega\rangle$
so that at $d = 1$ it is the $a$ of (\ref{eq:admissible-exponent}), and a \Levy{} measure $\nu$
(Lemma~\ref{lem:family-structure}). Sato's Gaussian covariance matrix, which enters as
$\tfrac12\langle z, A^{\mathrm S}z\rangle$ (Theorem~\ref{thm:cited-A39}), is written $A^{\mathrm S}$,
so that $A^{\mathrm S} = 2A$; the Gaussian coefficient $a_0$ of an isotropic origin datum
(Definition~\ref{def:isotropic-origin-data}) is in Sato's normalization, $A^{\mathrm S} = a_0I$. On
the line the folded \Levy{} measure is written $\varpi$
and the measure of (\ref{eq:generator-form}) $\kappa^{\mathrm C}$, the two related by
$\varpi(dx) = k(x)\,dx/x$ with $k(x) = \kappa^{\mathrm C}((x,\infty))$. The generator of a family
is $\mathcal{H}_t$, and $\mathcal{A}_d$ is the class of exponents admitted in dimension $d$. The
knots, the knot kernels and the ladder of a pyramid are defined after
Lemma~\ref{lem:family-dilations}.

\section{The facts of the literature}
\label{sec:cited-literature}

The proofs use four pieces of the literature: the \Levy--Khintchine theory of infinitely divisible
laws, in the book of Sato [\sscite{sato1999levy}]; Bernstein functions, in the book of Schilling, Song
and Vondra\v{c}ek [\sscite{schilling2012bernstein}]; Schoenberg's representation of \Polya{} frequency
densities, in Karlin's book [\sscite{karlin1968total}]; and one integral representation of the
modified Bessel function $K_\nu$ [\sscite{dlmf2026}]. Each fact is stated at its source's letter,
narrowed to what the proofs spend.

\subsection{Infinitely divisible laws}
\label{subsec:cited-sato}

A probability measure $\mu$ on $\RR^d$ is \emph{infinitely divisible} if for every $n$ it is the
$n$-fold convolution power of some probability measure. Write $D := \{x \in \RR^d : |x| \le 1\}$
for the closed unit ball.

% draft: the citations of [Sato, Thm. 8.1] at lines 255, 376, 378, 416, 428, 490, 794, 796, 798.
% CHANGED: the uniqueness clause (ii) is on p. 38, where the draft's lines 490, 794, 796, 798 cite
% p. 37 (VERBATIM-Sato.md, A39).
% CHANGED (fidelity review, 2026-11-16, F-T3C-4): the spender count is the \uses graph's (24 nodes,
% where "Sixteen" counted neither it nor the ledger's list), and the "not carried" normalization
% says it is the spending nodes' convention, not a difference of this node from the page.
% CHANGED (the theory paper's blueprint round, 2026-11-16, review of Section 2, NOTATION): Sato's
% covariance matrix A is written A^S, as in lem:family-structure, here and in A42, A44, A45 and A48,
% apart from this paper's Gaussian part A = A^S/2. A rename; the content is the page's.
\begin{theorem}[the \Levy--Khintchine representation]
\label{thm:cited-A39}
\lean{Pyramids.cited_A39_exists, Pyramids.cited_A39_unique, Pyramids.cited_A39_converse}
\leanok
\begin{enumerate}
\item[(i)] If $\mu$ is an infinitely divisible probability measure on $\RR^d$, then
\[
  \widehat\mu(z) = \exp\Bigl[-\tfrac12\langle z, A^{\mathrm S}z\rangle + i\langle\gamma, z\rangle
  + \int_{\RR^d}\bigl(e^{i\langle z, x\rangle} - 1 - i\langle z, x\rangle\mathbf{1}_D(x)\bigr)\,
  \nu(dx)\Bigr], \qquad z \in \RR^d,
\]
where $A^{\mathrm S}$ is a symmetric nonnegative-definite $d \times d$ matrix, $\gamma \in \RR^d$, and $\nu$
is a measure on $\RR^d$ with $\nu(\{0\}) = 0$ and $\int_{\RR^d}\min(|x|^2, 1)\,\nu(dx) < \infty$.
\item[(ii)] The representation of $\widehat\mu$ in (i) by $A^{\mathrm S}$, $\nu$ and $\gamma$ is unique.
\item[(iii)] Conversely, for every symmetric nonnegative-definite $d \times d$ matrix $A^{\mathrm S}$, every
measure $\nu$ with $\nu(\{0\}) = 0$ and $\int\min(|x|^2, 1)\,\nu(dx) < \infty$, and every
$\gamma \in \RR^d$, there is an infinitely divisible probability measure $\mu$ whose transform is
given by the formula of (i).
\end{enumerate}
The triple $(A^{\mathrm S}, \nu, \gamma)$ is the \emph{generating triplet} of $\mu$, $A^{\mathrm S}$ its
\emph{Gaussian covariance matrix} and $\nu$ its \emph{\Levy{} measure}.
% CHANGED (fidelity review, 2026-11-16, R71): the status line lists the declarations' hypotheses beside the node's (LINKAGE rule 6).
\statusA\quad\ledger{A39}\quad\emph{(Cited from [\sscite{sato1999levy}, Thm.~8.1, pp.~37--38]:
clause (i) with (8.1)--(8.2) at the foot of p.~37, (ii) and (iii) at the head of p.~38; the names
are Def.~8.2, p.~38. Admitted as a Lean axiom by the author on 2026-11-13. The three clauses are
already axioms of the affine module's formal development, for the same theorem at the same pages
(\texttt{AffineHemigroup.levy\_khintchine\_exists}, \texttt{levy\_khintchine\_unique} and
\texttt{levy\_khintchine\_converse}); the declarations named above are those axioms, and this
module declares none of its own for the entry.
\textbf{Assignment.} \ledger{A39} carries the three clauses. Twenty-four nodes name this one in
their \texttt{\textbackslash uses}, ten interface nodes of this chapter and fourteen nodes of the
later chapters, each naming the clauses it reads in its own annotation: the representation (i) in
\ref{lem:family-structure}; the uniqueness (ii) to identify data (the additivity of triplets under
convolution, the invariance of a \Levy{} measure under a symmetry of the law, the triplet after a
linear image); the existence (iii) for the laws of given data, as in \ref{lem:synthesis}. Not
carried: the reductions for a symmetric law, which the spending nodes prove in one line each (the
compensator integrates an odd function to $0$, so $\gamma = 0$ and the exponent is
$\int(1 - \cos\langle z, x\rangle)\,\nu(dx)$; a real or radial transform gives a symmetric or
rotation-invariant law, the transform determining the law); and the convention of the spending
nodes, which write a Gaussian part of an exponent as $\langle z, Az\rangle$, so that their $A$
is half the matrix $A^{\mathrm S}$ of (i) (the triplet $(2a, \tilde\nu, 0)$ in the proof of
Theorem~\ref{thm:cited-A9}). The statement above, and its declarations, keep Sato's
$\tfrac12\langle z, A^{\mathrm S}z\rangle$, Sato's matrix $A$ being written $A^{\mathrm S}$.
\textbf{Hypotheses} (the node's, each with what its declarations require): for (i), $\mu$ a probability
measure on $\RR^d$: same (an \texttt{IsProbabilityMeasure} instance, every $d$); $\mu$ infinitely
divisible: same; for (ii), the same two, and the two representations: same (\texttt{HasTriplet}
twice), the infinite divisibility being implied by them and so redundant; for (iii), $A^{\mathrm S}$ symmetric
nonnegative-definite, $\nu(\{0\}) = 0$, $\int\min(|x|^2, 1)\,d\nu < \infty$ and $\gamma \in \RR^d$:
same (the fields of \texttt{LevyTriplet}).)}
\end{theorem}

% draft: the citations of [Sato, Lemma 7.5, p. 32] at lines 378 and 794.
% CHANGED: narrowed to "no zero", the lemma's whole statement; "real" and "positive" are the
% spending nodes' steps (VERBATIM-Sato.md, A40).
\begin{theorem}[the transform has no zero]
\label{thm:cited-A40}
\lean{Pyramids.cited_A40}
\leanok
\uses{thm:cited-A39}
If $\mu$ is an infinitely divisible probability measure on $\RR^d$, then $\widehat\mu(z) \ne 0$ for
every $z \in \RR^d$.
% CHANGED (fidelity review, 2026-11-16, R71): the status line lists the declaration's hypotheses beside the node's (LINKAGE rule 6).
\statusA\quad\ledger{A40}\quad\emph{(Cited from [\sscite{sato1999levy}, Lemma~7.5, p.~32]. Admitted
as a Lean axiom by the author on 2026-11-13, and found to need none of its own: the declaration
named above is proved from clause (i) of Theorem~\ref{thm:cited-A39}, which writes $\widehat\mu$ as an
exponential, and an exponential has no zero. In the book the lemma comes first and is used to prove
Theorem~8.1; here that theorem is taken as given, so the order is reversed without circularity.
\textbf{Assignment.} \ledger{A40}
carries the statement. \ref{lem:family-structure} and \ref{prop:grid-pyramid-not-isotropic} spend
it. Not carried: that the transform of a symmetric law is real, and positive by continuity from
$\widehat\mu(0) = 1$, which the spending nodes prove.
\textbf{Hypotheses} (the node's, each with what its declaration requires): $\mu$ a probability measure
on $\RR^d$: same (an \texttt{IsProbabilityMeasure} instance, every $d$); $\mu$ infinitely
divisible: same; $z \in \RR^d$: same, every $z$.)}
\end{theorem}

% draft: the citations of [Sato, Thm. 9.3, p. 47] at lines 378 and 788.
\begin{theorem}[Khintchine's theorem]
\label{thm:cited-A41}
\lean{Pyramids.cited_A41}
\leanok
Let $\{Z_{nk} : k = 1, \dots, r_n;\ n = 1, 2, \dots\}$ be random variables on $\RR^d$ such that, for
each $n$, $Z_{n1}, \dots, Z_{nr_n}$ are independent, and, for every $\varepsilon > 0$,
\[
  \lim_{n \to \infty}\ \max_{1 \le k \le r_n}\Pr\bigl[|Z_{nk}| > \varepsilon\bigr] = 0
\]
(a \emph{null array}), and let $S_n := \sum_{k=1}^{r_n}Z_{nk}$ be its row sums. If, for some
$b_n \in \RR^d$, the law of $S_n - b_n$ converges to a probability measure $\mu$, then $\mu$ is
infinitely divisible.
% CHANGED (fidelity review, 2026-11-16, R71): the status line lists the declaration's hypotheses beside the node's (LINKAGE rule 6).
\statusA\quad\ledger{A41}\quad\emph{(Cited from [\sscite{sato1999levy}, Thm.~9.3, p.~47], with
Def.~9.2 and (9.1) on the same page. Admitted as a Lean axiom by the author on 2026-11-13. The
affine module's formal development already states the theorem as an axiom, for the laws of the
rows and with every row nonempty (\texttt{AffineHemigroup.khintchine\_null\_array}); the
declaration named above is proved from it, since the law of a sum of independent random variables
is the convolution of their laws, and a point mass at $0$ put in front of each row changes neither
the law of the row sum nor the maximum in the null condition. This module declares no axiom of its
own for the entry.
\textbf{Assignment.} \ledger{A41} carries the statement. \ref{lem:steps-infinitely-divisible}
spends it with $b_n = 0$ and row sums whose laws are all equal, and the nodes that read the
infinite divisibility of a step (\ref{lem:family-structure}, \ref{prop:no-bounded-step},
\ref{prop:grid-pyramid-not-isotropic}, \ref{prop:exact-masks}) spend it through that lemma. Not
carried: the construction of the null array from (I1) and (I2), and the choice of independent
realizations of the factor laws, which the lemma supplies.
\textbf{Hypotheses} (the node's, each with what its declaration requires): the array $Z_{nk}$ of random
variables on $\RR^d$: stronger, all rows live on one probability space, which the node leaves
open (a product space transports the general case); each row independent: same; $r_n$: same, any
$r_n \ge 0$, an empty row allowed; the null condition for every $\varepsilon > 0$: same, with the
supremum over the row; $b_n \in \RR^d$: same; the law of $S_n - b_n$ converges to a probability
measure $\mu$: same, against every bounded continuous function.)}
\end{theorem}

% draft: the citations of [Sato, Prop. 11.10, p. 65] at lines 410, 457, 467, 774, 796 and 798.
% CHANGED: stated for an infinitely divisible law; the page states it for a Levy process, and the
% law is the process's at time 1, by the correspondence stated on the same page
% (VERBATIM-Sato.md, A42).
\begin{theorem}[the linear image of an infinitely divisible law]
\label{thm:cited-A42}
\lean{Pyramids.cited_A42}
\leanok
\uses{thm:cited-A39}
Let $\mu$ be an infinitely divisible probability measure on $\RR^d$ with generating triplet
$(A^{\mathrm S}, \nu, \gamma)$, and let $U$ be an $n \times d$ matrix. Then the image $U_*\mu$ of $\mu$ under
$x \mapsto Ux$ is infinitely divisible on $\RR^n$, with generating triplet
\[
  A^{\mathrm S}_U = UA^{\mathrm S}U^{\tp}, \qquad \nu_U = \bigl[\nu U^{-1}\bigr]_{\RR^n \setminus \{0\}}, \qquad
  \gamma_U = U\gamma + \int Ux\,\bigl(\mathbf{1}_E(Ux) - \mathbf{1}_D(x)\bigr)\,\nu(dx),
\]
where $(\nu U^{-1})(B) = \nu(\{x : Ux \in B\})$, $[\,\cdot\,]_{\RR^n \setminus \{0\}}$ is the
restriction to $\RR^n \setminus \{0\}$, and $D$ and $E$ are the closed unit balls of $\RR^d$ and
$\RR^n$.
% CHANGED (2026-11-14, the author's decision): prop:no-lattice-kernel, regraded [T], no longer
% spends the entry; the orthogonal projections were its printed route's only use.
% CHANGED (fidelity review, 2026-11-15, R26/R50/R69): the status line lists the declarations' hypotheses beside the node's (LINKAGE rule 6).
% CHANGED (2026-11-23, review round 3 batch C, item 2): proved, Lean core, by the change of
% variables y = Ux in the formula of thm:cited-A39(i); no longer an admitted axiom. Statement and
% the declaration's type unchanged.
\statusA\quad\ledger{A42}\quad\emph{(Cited from [\sscite{sato1999levy}, Prop.~11.10, p.~65],
formulas (11.8)--(11.10), printed for a \Levy{} process $\{X_t\}$ and read for the law
$\mu = P_{X_1}$, through the correspondence of \Levy{} processes and infinitely divisible laws
stated on the same page. This citation is the node's grounding. Admitted as a Lean axiom by the
author on 2026-11-13; since review round 3 (2026-11-23) not admitted, and found to need no axiom:
the proof below derives the statement from the form of the transform in
Theorem~\ref{thm:cited-A39}(i), and its declaration prints Lean core.
\textbf{Assignment.} \ledger{A42} carries the statement, with the $\nu$-integrability of
$x \mapsto Ux\,(\mathbf{1}_E(Ux) - \mathbf{1}_D(x))$, which the formula for $\gamma_U$ presupposes
and the proof below shows; the typed statement asserts it. Five nodes spend it, for the matrices
$qI$, $q^nI$, an invertible $G^{-1}$, a rotation, a reflection $-I$ and a coordinate map:
\ref{lem:family-structure}, \ref{lem:lattice-infinitely-divisible},
\ref{prop:what-does-not-lift}, \ref{prop:what-does-not-lift-chain}
and \ref{prop:grid-pyramid-not-isotropic}, the last also reading $\gamma_U$ at $U = -I$, where the
correction term vanishes. Not carried: anything about the process beyond its law at time $1$.
\textbf{Hypotheses} (the node's, each with what its declarations require): $\mu$ a probability measure
on $\RR^d$: same (an \texttt{IsProbabilityMeasure} instance, every $d$); $\mu$ infinitely
divisible: same; $(A^{\mathrm S}, \nu, \gamma)$ a generating triplet of $\mu$: same (\texttt{HasTriplet} with
a \texttt{LevyTriplet}, whose fields are $A^{\mathrm S}$ symmetric nonnegative-definite, $\nu(\{0\}) = 0$ and
$\int\min(|x|^2, 1)\,d\nu < \infty$); $U$ an $n \times d$ matrix: same, any $n$.)}
\end{theorem}

% CHANGED (2026-11-23, review round 3 batch C, item 2): new proof environment; the statement was
% an admitted axiom until this round.
\begin{proof}
% CHANGED (2026-11-23, review round 3 batch C, register): the frequency on R^n is xi.
\leanok
\uses{thm:cited-A39}
Throughout, $\langle \xi, Ux\rangle = \langle U^{\tp}\xi, x\rangle$ for $x \in \RR^d$ and
$\xi \in \RR^n$, and $K$ is the operator norm of $U$, so that $|Ux| \le K|x|$. Write
$\psi_{A^{\mathrm S},\nu,\gamma}(z)$ for the bracket of Theorem~\ref{thm:cited-A39}(i), and
$g_z(x) = e^{i\langle z,x\rangle} - 1 - i\langle z,x\rangle\mathbf{1}_D(x)$ for its integrand on
$\RR^d$, with $\mathbf{1}_E$ in place of $\mathbf{1}_D$ on $\RR^n$.

\emph{Infinite divisibility.} Let $k \ge 1$ and let $\mu_k$ be a probability measure with
$\mu = \mu_k^{*k}$. The transform of $(U_*\mu_k)^{*k}$ at $\xi$ is
$\widehat{\mu_k}(U^{\tp}\xi)^k = \widehat\mu(U^{\tp}\xi)$, which is the transform of $U_*\mu$ at $\xi$; so
$U_*\mu = (U_*\mu_k)^{*k}$, and $U_*\mu$ is infinitely divisible.

\emph{The correction is integrable.} Put $c(x) = Ux\,(\mathbf{1}_E(Ux) - \mathbf{1}_D(x))$. It
vanishes unless exactly one of $|x| \le 1$ and $|Ux| \le 1$ holds. If $|x| > 1 \ge |Ux|$, then
$|c(x)| = |Ux| \le 1 = \min(|x|^2, 1)$. If $|x| \le 1 < |Ux|$, then $1 < |Ux| \le K|x|$, so
$|c(x)| \le K|x| \le (K|x|)^3 \le K^3|x|^2 = K^3\min(|x|^2, 1)$. Hence
$|c(x)| \le (K^3 + 1)\min(|x|^2, 1)$ for every $x$, and $c$ is $\nu$-integrable.

% CHANGED (2026-11-23, review round 3, batch C typesetting): wording.
% CHANGED (2026-11-23, review round 3, batch C typesetting): wording.
\emph{The triplet.} The matrix $UA^{\mathrm S}U^{\tp}$ is symmetric. It is nonnegative-definite, as
$\langle \xi, UA^{\mathrm S}U^{\tp}\xi\rangle = \langle U^{\tp}\xi, A^{\mathrm S}U^{\tp}\xi\rangle$. The measure
$\nu_U$ has no mass at $0$, and since $\min(|Ux|^2, 1) \le (K^2 + 1)\min(|x|^2, 1)$,
\[
  \int \min(|y|^2, 1)\,\nu_U(dy) \le \int \min(|Ux|^2, 1)\,\nu(dx)
  \le (K^2 + 1)\int \min(|x|^2, 1)\,\nu(dx) < \infty .
\]
So $(A^{\mathrm S}_U, \nu_U, \gamma_U)$ is a generating triplet on $\RR^n$ in the sense of
Theorem~\ref{thm:cited-A39}(i).

\emph{The transform.} Fix $\xi \in \RR^n$. The integrand $e^{i\langle \xi,y\rangle} - 1 - i\langle \xi,y\rangle\mathbf{1}_E(y)$
vanishes at $y = 0$, so its integral against $\nu_U$ is its integral against the image of $\nu$,
which is $\int (e^{i\langle \xi,Ux\rangle} - 1 - i\langle \xi,Ux\rangle\mathbf{1}_E(Ux))\,\nu(dx)$. Since
$\langle \xi, Ux\rangle = \langle U^{\tp}\xi, x\rangle$, the integrand here is
$g_{U^{\tp}\xi}(x) - i\langle \xi, c(x)\rangle$, a difference of two $\nu$-integrable functions, so the
integral is $\int g_{U^{\tp}\xi}\,d\nu - i\langle \xi, \int c\,d\nu\rangle$. With
$\langle \xi, A^{\mathrm S}_U \xi\rangle = \langle U^{\tp}\xi, A^{\mathrm S}U^{\tp}\xi\rangle$ and
$\langle \gamma_U, \xi\rangle = \langle \gamma, U^{\tp}\xi\rangle + \langle \xi, \int c\,d\nu\rangle$, the
two terms in $\langle \xi, \int c\,d\nu\rangle$ cancel, and
\[
  \psi_{A^{\mathrm S}_U,\nu_U,\gamma_U}(\xi) = \psi_{A^{\mathrm S},\nu,\gamma}(U^{\tp}\xi).
\]
The transform of $U_*\mu$ at $\xi$ is $\widehat\mu(U^{\tp}\xi) = \exp\psi_{A^{\mathrm S},\nu,\gamma}(U^{\tp}\xi)$,
so $(A^{\mathrm S}_U, \nu_U, \gamma_U)$ is a generating triplet of $U_*\mu$, the only one by
Theorem~\ref{thm:cited-A39}(ii).

Every step is machine-checked. The formal statement, that $U_*\mu$ has this triplet, reads only
the form of the transform in Theorem~\ref{thm:cited-A39}(i) and no admitted fact. The
uniqueness in the last sentence is \ledger{A39}(ii).
\end{proof}

% draft: the citations of [Sato, Prop. 15.5, p. 93] at lines 134, 402 and 457.
\begin{theorem}[self-decomposable laws are infinitely divisible]
\label{thm:cited-A43}
\lean{Pyramids.cited_A43_sd, Pyramids.cited_A43_semi, Pyramids.cited_A43_span}
\leanok
\uses{def:self-decomposable, thm:cited-A40}
If $\mu$ is self-decomposable, then $\mu$ is infinitely divisible and, for every $b > 1$, the
measure $\varrho_b$ of Definition~\ref{def:self-decomposable} is uniquely determined and infinitely
divisible. If $\mu$ is semi-selfdecomposable, then $\mu$ is infinitely divisible and $\varrho_b$ is
uniquely determined. So a self-decomposable law is semi-selfdecomposable with every $b > 1$ as a
span.
% CHANGED (fidelity review, 2026-11-15, R26/R50/R69): the status line lists the declarations' hypotheses beside the node's (LINKAGE rule 6).
\statusA\quad\ledger{A43}\quad\emph{(Cited from [\sscite{sato1999levy}, Prop.~15.5, p.~93], with the
sentence after it on the same page for the last clause. Admitted as a Lean axiom by the author on
2026-11-14, in three parts. That a self-decomposable law is infinitely divisible is already an axiom
of the affine module's formal development, for the same sentence of the same page
(\texttt{AffineHemigroup.selfDecomposable\_infinitelyDivisible}). That its residuals are
infinitely divisible, and that a semi-selfdecomposable law is, are declared as two axioms of this
module. The two uniqueness clauses are proved by the page's own argument: an infinitely divisible
transform has no zero (Theorem~\ref{thm:cited-A40}), so $\widehat\varrho_b = \widehat\mu/\widehat\mu(b^{-1}\,\cdot)$,
and the transform determines the law. The last sentence follows from the first clause. The
declarations named above are theorems from these. \textbf{Assignment.} \ledger{A43} carries the
whole statement.
\ref{prop:what-does-not-lift-chain} spends the last clause and the infinite divisibility;
\ref{prop:strata-cylinder-coordinates-q-cin}(3) spends the infinite divisibility only; chapter~3
states the proposition in Section~\ref{sec:pyramid-strata}. Not carried: the
characterization of semi-selfdecomposability by the \Levy{} measure, which is
Theorem~\ref{thm:cited-A44}.
\textbf{Hypotheses} (the node's, each with what its declarations require): $\mu$ a probability measure
on $\RR^d$: same (an \texttt{IsProbabilityMeasure} instance, every $d$); $\mu$ self-decomposable:
same (\texttt{IsSelfDecomposableSato}, Definition~\ref{def:self-decomposable} at Sato's letter),
in \texttt{cited\_A43\_sd} and \texttt{cited\_A43\_span}; $b > 1$: same; $\varrho_b$ a probability
measure with the identity of Definition~\ref{def:self-decomposable}: same
(\texttt{IsSDResidual}); $\mu$ semi-selfdecomposable: same (\texttt{IsSemiSelfDecomposable}, some
$b > 1$ with an infinitely divisible residual), in \texttt{cited\_A43\_semi}, whose uniqueness is
required at every span $b$.)}
\end{theorem}

% draft: the citations of [Sato, Thm. 15.8, p. 94] at lines 134, 402 and 457.
\begin{theorem}[semi-selfdecomposability by the \Levy{} measure]
\label{thm:cited-A44}
\lean{Pyramids.cited_A44}
% CHANGED (wave 3d, 2026-11-14): proved from A39 and A42 by the page's own argument; statement and
% the declaration's type unchanged. The entry is not admitted and needs no axiom of its own.
\leanok
\uses{def:self-decomposable, thm:cited-A39, thm:cited-A42}
Fix $b > 1$. Let $\mu$ be an infinitely divisible probability measure on $\RR^d$ with generating
triplet $(A^{\mathrm S}, \nu, \gamma)$. Then $\mu$ is semi-selfdecomposable with $b$ as a span if and only if
\[
  D_b\nu \ge \nu ,
\]
where $(D_b\nu)(B) = \nu(b^{-1}B)$ is the image of $\nu$ under $x \mapsto bx$. Semi-selfdecomposability
imposes no restriction on $A^{\mathrm S}$ and $\gamma$.
% CHANGED (fourth blind pass, 2026-11-15): the derivation from A39 and A42 moves out of the
% annotation into the proof below, completed (the sum of two triplets, the absolute continuity and
% the difference measure).
% CHANGED (fidelity review, 2026-11-16, R71): the status line lists the declaration's hypotheses beside the node's (LINKAGE rule 6).
\statusA\quad\ledger{A44}\quad\emph{(Cited from [\sscite{sato1999levy}, Thm.~15.8, p.~94], formula
(15.6), where the image is written $T_b\nu$ and defined at (13.8), p.~75. This citation is the
node's grounding. Not admitted, and found to need no axiom of its own: the proof below derives
the statement from Theorem~\ref{thm:cited-A39} and Theorem~\ref{thm:cited-A42} by the argument of Sato's proof.
\textbf{Assignment.} \ledger{A44} carries the
equivalence and its last sentence. \ref{prop:lindeberg-stratum-one}(3) spends the equivalence
and \ref{prop:stratum-one-example} spends it. Not carried: the
equivalent form (15.7) on the shells $\{b^n < |x| \le b^{n+1}\}$, which no node reads.
\textbf{Hypotheses} (the node's, each with what its declaration requires): $b > 1$: same; $\mu$ a
probability measure on $\RR^d$: same (an \texttt{IsProbabilityMeasure} instance, every $d$); $\mu$
infinitely divisible: same; $(A^{\mathrm S}, \nu, \gamma)$ a generating triplet of $\mu$: same
(\texttt{HasTriplet} with a \texttt{LevyTriplet}); the inequality $D_b\nu \ge \nu$: same, an
inequality of measures, that is on every measurable set.)}
\end{theorem}

% CHANGED (fourth blind pass, 2026-11-15): new proof environment, the derivation formerly sketched
% in the annotation, with its missing steps written out.
% CHANGED (the theory paper, 2026-11-16): the closing sentence on checking and cited facts shortened; no step changed.
% CHANGED (the theory paper's blueprint round, 2026-11-16, review of Section 2): Sato's matrices are
% written A^S, as in the statement, and the residual varrho_b, as in def:self-decomposable. Renames;
% no step changed.
\begin{proof}
\leanok
\uses{def:self-decomposable, thm:cited-A39, thm:cited-A42}
For a generating triplet $(A^{\mathrm S}, \nu, \gamma)$ write $\psi_{A^{\mathrm S},\nu,\gamma}(z)$ for the bracket of
Theorem~\ref{thm:cited-A39}(i), so that the law with this triplet has the transform
$\exp\psi_{A^{\mathrm S},\nu,\gamma}$. If $(A^{\mathrm S}_1, \nu_1, \gamma_1)$ and
$(A^{\mathrm S}_2, \nu_2, \gamma_2)$ are triplets, so is
$(A^{\mathrm S}_1 + A^{\mathrm S}_2, \nu_1 + \nu_2, \gamma_1 + \gamma_2)$: the sum of the matrices is symmetric and
nonnegative-definite, and $\nu_1 + \nu_2$ has no mass at $0$ and integrates $\min(|x|^2, 1)$. The
integrand of the bracket is integrable for $\nu_1$ and for $\nu_2$, so its integral against
$\nu_1 + \nu_2$ is the sum of the two integrals, and
\[
  \psi_{A^{\mathrm S}_1 + A^{\mathrm S}_2,\,\nu_1 + \nu_2,\,\gamma_1 + \gamma_2}
  = \psi_{A^{\mathrm S}_1,\nu_1,\gamma_1} + \psi_{A^{\mathrm S}_2,\nu_2,\gamma_2}.
\]

Let $\mu_b$ be the image of $\mu$ under $x \mapsto b^{-1}x$, so that
$\widehat\mu_b(z) = \widehat\mu(b^{-1}z)$. By Theorem~\ref{thm:cited-A42} with $n = d$ and $U = b^{-1}I$,
the law $\mu_b$ is infinitely divisible with a triplet $(b^{-2}A^{\mathrm S}, \nu_b, \gamma')$, where $\nu_b$
is the image of $\nu$ under $x \mapsto b^{-1}x$ (the restriction to $\RR^d \setminus \{0\}$ changes
nothing, as $\nu(\{0\}) = 0$) and $\gamma' \in \RR^d$. Since $\nu_b(B) = \nu(bB)$ and
$(D_b\nu)(B) = \nu(b^{-1}B)$, the inequality $D_b\nu \ge \nu$ on all Borel sets is the inequality
$\nu \ge \nu_b$, read at the set $bB$.

Let $\mu$ be semi-selfdecomposable with span $b$. By Definition~\ref{def:self-decomposable} its
residual $\varrho_b$ is infinitely divisible, so it has a triplet
$(A^{\mathrm S}_\varrho, \nu_\varrho, \gamma_\varrho)$ by
Theorem~\ref{thm:cited-A39}(i). Then
$\widehat\mu = \widehat\mu_b\,\widehat\varrho_b = \exp(\psi_{b^{-2}A^{\mathrm S},\nu_b,\gamma'} +
\psi_{A^{\mathrm S}_\varrho,\nu_\varrho,\gamma_\varrho})$, and by the first paragraph this is the transform given by the
summed triplet $(b^{-2}A^{\mathrm S} + A^{\mathrm S}_\varrho, \nu_b + \nu_\varrho, \gamma' + \gamma_\varrho)$. By
Theorem~\ref{thm:cited-A39}(ii) the summed triplet is $(A^{\mathrm S}, \nu, \gamma)$. So
$\nu = \nu_b + \nu_\varrho \ge \nu_b$, which is $D_b\nu \ge \nu$.

Conversely, let $\nu_b \le \nu$. The measure $\nu$ is $\sigma$-finite: it is finite on
$\{|x| > 1\}$ and on each shell $\{1/(n+1) < |x| \le 1/n\}$, $n \ge 1$, by the integrability of
$\min(|x|^2, 1)$, and $\nu(\{0\}) = 0$. From $\nu_b \le \nu$, every $\nu$-null set is
$\nu_b$-null, so by the Radon--Nikodym theorem $\nu_b = f\nu$ for a measurable $f \ge 0$, and
$f \le 1$ $\nu$-almost everywhere, again because $\nu_b \le \nu$. Put $\nu_\varrho := (1 - f)\nu$.
Then $\nu_b + \nu_\varrho = \nu$ and $\nu_\varrho \le \nu$, so $\nu_\varrho$ has no mass at $0$ and
integrates $\min(|x|^2, 1)$. The matrix
$(1 - b^{-2})A^{\mathrm S}$ is symmetric and nonnegative-definite because $b > 1$. By
Theorem~\ref{thm:cited-A39}(iii) there is an infinitely divisible law $\varrho_b$ with the triplet
$((1 - b^{-2})A^{\mathrm S}, \nu_\varrho, \gamma - \gamma')$. By the first paragraph
$\widehat\mu(b^{-1}z)\,\widehat\varrho_b(z) = \exp\psi_{A^{\mathrm S},\nu_b + \nu_\varrho,\gamma}(z) =
\widehat\mu(z)$, so $\mu$ is semi-selfdecomposable with span $b$ and residual $\varrho_b$.

The criterion involves $\nu$ alone, which is the last sentence of the statement. Every step is
machine-checked, and the cited facts are \ledger{A39}(i)--(iii) and \ledger{A42}.
\end{proof}

% draft: the citations of [Sato, Thm. 15.10 and Cor. 15.11, p. 95] at lines 152, 158, 436 and 500.
% CHANGED: narrowed to Thm. 15.10; Cor. 15.11, the line case, is named in the annotation (no node
% spends it: thm:lattice-line says so). The draft's consequences, that the Levy measure gives
% countable sets and spheres about 0 mass 0, are the spending nodes' one-line steps
% (VERBATIM-Sato.md, A45).
\begin{theorem}[the polar form of a self-decomposable law]
\label{thm:cited-A45}
\lean{Pyramids.cited_A45}
\leanok
\uses{def:self-decomposable, thm:cited-A39}
Let $\mu$ be an infinitely divisible probability measure on $\RR^d$ with generating triplet
$(A^{\mathrm S}, \nu, \gamma)$, and let $S = \{\xi \in \RR^d : |\xi| = 1\}$. Then $\mu$ is self-decomposable if
and only if
\[
  \nu(B) = \int_S\lambda(d\xi)\int_0^\infty\mathbf{1}_B(r\xi)\,k_\xi(r)\,\frac{dr}{r}
\]
for every Borel set $B$, with a finite measure $\lambda$ on $S$ and a nonnegative function
$k_\xi(r)$, measurable in $\xi \in S$ and nonincreasing in $r > 0$. Self-decomposability imposes no
restriction on $A^{\mathrm S}$ and $\gamma$.
% CHANGED (fidelity review, 2026-11-16, R71): the status line lists the declaration's hypotheses beside the node's (LINKAGE rule 6).
\statusA\quad\ledger{A45}\quad\emph{(Cited from [\sscite{sato1999levy}, Thm.~15.10, p.~95], formula
(15.8); Sato's ``decreasing'' is meant in the weak sense. Admitted as a Lean axiom by the author on
2026-11-14. The theorem is already an axiom of the affine module's formal development, for the same
page (\texttt{AffineHemigroup.sato\_polar\_criterion}); the declaration named above is that axiom,
and this module declares none of its own for the entry. \textbf{Assignment.} \ledger{A45} carries the statement.
\ref{prop:stratum-two-example} and \ref{prop:ladder-knot-kernel}(4) spend the ``only if''. Not
carried: the consequences that such a $\nu$ gives every countable set and every sphere about $0$
mass $0$ (for each $\xi$ the set of $r$ concerned is countable or a point, hence $dr/r$-null), which
the spending nodes prove; and the case $d = 1$ written with a density, Sato's Cor.~15.11 on the same
page, which no node spends.
\textbf{Hypotheses} (the node's, each with what its declaration requires): $\mu$ a probability measure
on $\RR^d$: same (an \texttt{IsProbabilityMeasure} instance, every $d$); $\mu$ infinitely
divisible: same; $(A^{\mathrm S}, \nu, \gamma)$ a generating triplet of $\mu$: same (\texttt{HasTriplet}); on
the right side, $\lambda$ a finite measure on $S$: same; $k_\xi(r) \ge 0$: same, for $r > 0$;
$k_\xi(r)$ measurable in $\xi$: same, for each $r > 0$, the joint measurability being implied;
$k_\xi$ nonincreasing in $r > 0$: same, read in the weak sense (\texttt{AntitoneOn}).)}
\end{theorem}

% draft: the citations of [Sato, Cor. 24.4, p. 149] at lines 219, 279, 508 and 815.
\begin{theorem}[unbounded support]
\label{thm:cited-A46}
\lean{Pyramids.cited_A46}
\leanok
If $\mu$ is an infinitely divisible probability measure on $\RR^d$ and is not a point mass, then
the support of $\mu$ is unbounded.
% CHANGED (fidelity review, 2026-11-15, R26/R50/R69): the status line lists the declarations' hypotheses beside the node's (LINKAGE rule 6).
\statusA\quad\ledger{A46}\quad\emph{(Cited from [\sscite{sato1999levy}, Cor.~24.4, p.~149], where a
point mass is called a $\delta$-distribution. Admitted as a Lean axiom by the author on
2026-11-13, at the statement typed in the skeleton, where an unbounded support is positive mass
outside every ball. \textbf{Assignment.} \ledger{A46} carries the statement. \ref{prop:no-bounded-step}
spends it, and \ref{prop:exact-masks}(1) through that node. Not carried: anything about the shape of the support, Sato's
Thm.~24.5 and Cor.~24.6, which is Theorem~\ref{thm:cited-A48} in the part used.
\textbf{Hypotheses} (the node's, each with what its declarations require): $\mu$ a probability measure
on $\RR^d$: same (an \texttt{IsProbabilityMeasure} instance, every $d$); $\mu$ infinitely
divisible: same; $\mu$ not a point mass: same, $\mu \ne \delta_x$ for every $x$.)}
\end{theorem}

% draft: the citations of [Sato, Cor. 24.6, p. 149] at lines 279 and 774.
% CHANGED: narrowed to the "only if", on R as printed; the finiteness of the Levy measure is printed
% in the proof on the same page, not in the statement, and the drift clause is dropped
% (VERBATIM-Sato.md, A48).
% CHANGED (2026-11-23, review round 3 batch C, item 1): the admitted part is the corollary's
% statement only, A = 0 and the support of nu in Z; the finiteness of nu, until now admitted from
% the proof on the same page, follows from the support clause and is proved below.
\begin{theorem}[infinitely divisible laws carried by $\ZZ$]
\label{thm:cited-A48}
\lean{Pyramids.cited_A48, Pyramids.cited_A48_finite}
\leanok
\uses{thm:cited-A39}
Let $\mu$ be an infinitely divisible probability measure on $\RR$ with generating triplet
$(A^{\mathrm S}, \nu, \gamma)$. If the support of $\mu$ is contained in $\ZZ$, then $A^{\mathrm S} = 0$ and the
support of $\nu$ is contained in $\ZZ$. In particular $\nu(\RR) < \infty$.
% CHANGED (fidelity review, 2026-11-15, R26/R50/R69): the status line lists the declarations' hypotheses beside the node's (LINKAGE rule 6).
\statusA\quad\ledger{A48}\quad\emph{(Cited from [\sscite{sato1999levy}, Cor.~24.6, p.~149], its
``only if'' without the drift clause: the first sentence of the conclusion is what the corollary
states. Admitted as a Lean axiom by the author on 2026-11-13, at the statement typed in the
skeleton, where a support contained in $\ZZ$ is a null complement of $\ZZ$; since review round 3
(2026-11-23) the axiom \texttt{cited\_A48} is that sentence and nothing more. The finiteness of
$\nu$, which Sato shows in the proof on the same page, is proved below from the support clause.
\textbf{Assignment.} \ledger{A48} carries the first sentence of the conclusion.
\ref{lem:lattice-infinitely-divisible}
spends the node, on the coordinate images of a law carried by a lattice. Not carried: the converse; the
corollary's condition that the drift lies in $\ZZ$; any version on $\RR^d$, the page being on
$\RR$; and the finiteness of $\nu$, which is proved.
\textbf{Hypotheses} (the node's, each with what its declarations require): $\mu$ a probability measure
on $\RR$: same, typed on $\RR^1$ (\texttt{Rd 1}) through $x \mapsto x_0$; $\mu$ infinitely
divisible: same; $(A^{\mathrm S}, \nu, \gamma)$ a generating triplet of $\mu$: same (\texttt{HasTriplet});
the support of $\mu$ contained in $\ZZ$: same, as $\mu(\RR \setminus \ZZ) = 0$, equivalent since
$\ZZ$ is closed.)}
\end{theorem}

% CHANGED (2026-11-23, review round 3 batch C, item 1): new proof environment, for the finiteness
% of nu, formerly admitted.
\begin{proof}
\leanok
\uses{thm:cited-A39}
The first sentence is cited. For the second: $\nu(\{0\}) = 0$ and $\nu$ is carried by $\ZZ$, so it
is carried by $\ZZ \setminus \{0\}$, where $|x| \ge 1$ and so $\min(|x|^2, 1) = 1$. Hence
$\nu(\RR) \le \int \min(|x|^2, 1)\,\nu(dx) < \infty$ by Theorem~\ref{thm:cited-A39}(i). Every step
is machine-checked, and the cited fact is \ledger{A48}, for the first sentence.
\end{proof}

% draft: the citations of [Sato, Thm. 25.3, p. 159] at lines 285, 496 and 823.
% CHANGED: the hypothesis "locally bounded" of the page is stated, which the draft omits; the
% statement is for the law, the page's process read at time 1 (VERBATIM-Sato.md, A49).
\begin{theorem}[the $g$-moment theorem]
\label{thm:cited-A49}
\lean{Pyramids.cited_A49}
\leanok
\uses{thm:cited-A39}
A function $g$ on $\RR^d$ is \emph{submultiplicative} if it is nonnegative and there is a constant
$a > 0$ with $g(x + y) \le a\,g(x)g(y)$ for all $x, y \in \RR^d$, and \emph{locally bounded} if it is
bounded on every compact set. Let $g$ be a submultiplicative, locally bounded, measurable function on
$\RR^d$, and let $\mu$ be an infinitely divisible probability measure on $\RR^d$ with \Levy{} measure
$\nu$. Then $\int g\,d\mu < \infty$ if and only if $\int_{\{|x| > 1\}}g\,d\nu < \infty$.
% CHANGED (fidelity review, 2026-11-15, R26/R50/R69): the status line lists the declarations' hypotheses beside the node's (LINKAGE rule 6).
\statusA\quad\ledger{A49}\quad\emph{(Cited from [\sscite{sato1999levy}, Thm.~25.3, p.~159], with
Defs.~25.1--25.2 on the same page, printed for a \Levy{} process $\{X_t\}$, finite $g$-moment of
$X_t$ for every $t > 0$; the page says that finiteness of the $g$-moment ``is not a time dependent
distributional property'', so it is read at $t = 1$ for the law $P_{X_1} = \mu$. Admitted as a
Lean axiom by the author on 2026-11-14, at the statement typed in the skeleton, both directions.
\textbf{Assignment.} \ledger{A49} carries the
statement. \ref{prop:ladder-knot-kernel}(2) spends it, and \ref{prop:exact-masks} through that node, with
$g(x) = \log(|x| \vee e)$, which is submultiplicative by Theorem~\ref{thm:cited-A50} and locally bounded,
being continuous. Not carried: the comparison of $g$ with $\log_+|x|$, which the spending nodes
make.
\textbf{Hypotheses} (the node's, each with what its declarations require): $g$ submultiplicative:
same (\texttt{IsSubmultiplicative}: $g \ge 0$ and some $a > 0$); $g$ locally bounded: same, a bound
of $|g|$ on every compact set; $g$ measurable: same; $\mu$ a probability measure on $\RR^d$: same
(an \texttt{IsProbabilityMeasure} instance, every $d$); $\mu$ infinitely divisible: same; $\nu$
the \Levy{} measure of $\mu$: same, the measure of a triplet with \texttt{HasTriplet}.)}
\end{theorem}

% draft: the citations of [Sato, Prop. 25.4(iii), p. 159] at lines 285, 496 and 823.
% CHANGED: narrowed to the one function of the list that the nodes use.
% CHANGED (fidelity review, 2026-11-15, R66): the status line says that the edge to thm:cited-A49
% is the definition of "submultiplicative", which that node states, and no ingredient of the proof.
\begin{theorem}[a logarithmic submultiplicative function]
\label{thm:cited-A50}
\lean{Pyramids.cited_A50}
\leanok
\uses{thm:cited-A49}
The function $g(x) = \log(|x| \vee e)$ on $\RR^d$ is submultiplicative.
% CHANGED (fidelity review, 2026-11-16, R71): the status line lists the declaration's hypotheses beside the node's (LINKAGE rule 6).
\statusA\quad\ledger{A50}\quad\emph{(Cited from [\sscite{sato1999levy}, Prop.~25.4(iii), p.~159],
where $\log(|x| \vee e)$ is one of the functions listed. Admitted as a Lean axiom by the author on
2026-11-14, and found to need none: the declaration named above is proved with no cited fact, with
the constant $a = 2$. Put $u = |x| \vee e$ and $v = |y| \vee e$; both exceed $2$, so
$|x + y| \le u + v \le uv$ and $e \le uv$, whence $g(x + y) \le \log u + \log v \le 2\log u\log v$,
both logarithms being at least $1$. The node uses Theorem~\ref{thm:cited-A49} for the definition of a
submultiplicative function only; the proof reads nothing from it. \textbf{Assignment.} \ledger{A50}
carries the statement.
\ref{prop:ladder-knot-kernel}(2) spends it, and \ref{prop:exact-masks} through that node. Not
carried: the rest of the
proposition's list and its clauses (i)--(ii); and local boundedness, which the proposition does not
state and which holds since $g$ is continuous.
\textbf{Hypotheses} (the node's, each with what its declaration requires): none beyond the dimension,
every $d$ (\texttt{IsSubmultiplicative}: $g \ge 0$ and some $a > 0$ with
$g(x + y) \le a\,g(x)\,g(y)$, proved with $a = 2$).)}
\end{theorem}

\subsection{Bernstein functions}
\label{subsec:cited-bernstein}

A function $f : (0,\infty) \to \RR$ is \emph{completely monotone} if it is smooth
and $(-1)^nf^{(n)} \ge 0$ for every $n \ge 0$. A \emph{Bernstein function} is a
function $f \ge 0$ on $(0,\infty)$ whose derivative is completely monotone. These are exactly the
functions
\[
  f(\lambda) = a + b\lambda + \int_{(0,\infty)}(1 - e^{-\lambda t})\,\mu(dt), \qquad a, b \ge 0,
  \quad \int_{(0,\infty)}\min(1, t)\,\mu(dt) < \infty,
\]
and the triplet $(a, b, \mu)$ is determined by $f$ [\sscite{schilling2012bernstein}, Thm.~3.2,
p.~21]. The measure $\mu$ is the \emph{\Levy{} measure} of $f$.

% draft: the paragraph before Proposition 4.2 (section 4.1), line 561, which defines the complete
% Bernstein functions in words.
\begin{definition}[complete Bernstein functions]
\label{def:complete-bernstein}
\lean{Pyramids.IsCompletelyMonotone, Pyramids.IsBernsteinFunction, Pyramids.IsCompleteBernstein}\leanok
A Bernstein function $f$ is a \emph{complete Bernstein function}
[\sscite{schilling2012bernstein}, Def.~6.1, p.~69] if its \Levy{} measure has a
completely monotone density $m$ with respect to Lebesgue measure:
\[
  f(\lambda) = a + b\lambda + \int_{(0,\infty)}(1 - e^{-\lambda t})\,m(t)\,dt .
\]
\statusT\quad\emph{(A definition, [\sscite{schilling2012bernstein}, Def.~6.1, p.~69], formula (6.1),
which reads the \Levy{} measure in the representation (3.2) of the book's Thm.~3.2, p.~21. It
asserts nothing and cites no ledger entry.)}
\end{definition}

% draft: the same paragraph, line 561.
\begin{definition}[Thorin--Bernstein functions]
\label{def:thorin-bernstein}
\lean{Pyramids.IsThorinBernstein}\leanok
\uses{def:complete-bernstein}
A Bernstein function $f$ is a \emph{Thorin--Bernstein function}
[\sscite{schilling2012bernstein}, Def.~8.1, p.~109] if its \Levy{} measure has a density
$m$ such that $t \mapsto t\,m(t)$ is completely monotone, that is
\[
  f(\lambda) = a + b\lambda + \int_{(0,\infty)}(1 - e^{-\lambda t})\,m(t)\,dt,
\]
where $a, b \ge 0$, $\int_{(0,\infty)}\min(1, t)\,m(t)\,dt < \infty$ and $t \mapsto t\,m(t)$ is
completely monotone.
\statusT\quad\emph{(A definition, [\sscite{schilling2012bernstein}, Def.~8.1, p.~109], formula (8.1).
The same page observes that $m(t) = t^{-1}\cdot t\,m(t)$ is then completely monotone, so a
Thorin--Bernstein function is a complete Bernstein function. It asserts nothing and cites no ledger
entry.)}
\end{definition}

% draft: the citation of [Schilling-Song-Vondracek, Thm. 1.4, p. 3] at line 588.
% CHANGED: narrowed to the existence of the measure; its uniqueness and the converse of the page are
% not used (VERBATIM-literature.md, A55).
\begin{theorem}[Bernstein's theorem]
\label{thm:cited-A55}
\lean{Pyramids.cited_A55}
\notready
Let $f : (0,\infty) \to \RR$ be completely monotone. Then there is a measure $\mu$ on $[0,\infty)$
with
\[
  f(\lambda) = \int_{[0,\infty)}e^{-\lambda t}\,\mu(dt) \qquad\text{for every } \lambda > 0 .
\]
\statusA\quad\ledger{A55}\quad\emph{(Cited from [\sscite{schilling2012bernstein}, Thm.~1.4, p.~3].
The measure is positive, a completely monotone function being nonnegative. Its statement is typed in the skeleton
(\texttt{PyramidsSkeleton.Ch02CitedFacts}) and is not an axiom; admitting it is the author's decision at phase~5. \textbf{Assignment.} \ledger{A55} carries the statement.
% CHANGED (wave 3b, 2026-11-15): prop:thorin-in-dimension-d(1)-(2) and prop:matern-increments(3)
% no longer spend A55; they read lem:cm-log-convex, lem:cm-exponential-rate and
% lem:cm-uncancelled-atom, proved without it. The entry has no spender. The text read:
% "prop:matern-increments(3) and prop:thorin-in-dimension-d(1)-(2) spend it."
No node spends it: \ref{prop:thorin-in-dimension-d}(1)--(2) and \ref{prop:matern-increments}(3),
which read it, are proved from Lemmas~\ref{lem:cm-log-convex}, \ref{lem:cm-exponential-rate} and
\ref{lem:cm-uncancelled-atom}, without a Laplace representation. Not carried:
the uniqueness of $\mu$ and the converse, both on the page and neither used.)}
\end{theorem}

% draft: the paragraph before Proposition 4.2 (section 4.1), line 561: "e^{-f(|omega|^2)} is then
% the transform of an infinitely divisible law in every dimension, by subordination".
% CHANGED: the hypothesis f(0+) = 0 is stated; without it the measures have mass e^{-f(0+)} < 1
% (VERBATIM-literature.md, A57).
\begin{theorem}[a Bernstein function of $|\omega|^2$]
\label{thm:cited-A57}
\lean{Pyramids.cited_A57}
\notready
% CHANGED (skeleton review, m6): the value at 0, which f(|omega|^2) reads at omega = 0, is set
% to f(0+); the page reads it so, as the total mass exp(-t f(0)) of (13.21) (VERBATIM-literature.md, A57).
Let $f$ be a Bernstein function with $f(0{+}) = 0$, and set $f(0) := f(0{+}) = 0$. Then, for
every $d \ge 1$, the function
$\omega \mapsto f(|\omega|^2)$ is continuous and negative definite on $\RR^d$, and
$e^{-f(|\omega|^2)}$ is the transform of an infinitely divisible probability measure on $\RR^d$.
\statusA\quad\ledger{A57}\quad\emph{(Cited from [\sscite{schilling2012bernstein}, Thm.~13.14, p.~212],
the direction proved there by Example~13.13 on the same page: $f(|\xi|^2)$ is continuous and
negative definite for all $d$, and $e^{-tf(|\xi|^2)}$ is for each $t > 0$ the transform of a
rotation-invariant measure $\nu^t_d$ of mass $e^{-tf(0)}$, (13.21). With $f(0{+}) = 0$ these are
probability measures forming a convolution semigroup, so $\nu^1_d = (\nu^{1/n}_d)^{*n}$ is
infinitely divisible. The mechanism is subordination, [\sscite{sato1999levy}, Thm.~30.1,
pp.~197--198]: Brownian motion on $\RR^d$ with transform $e^{-|\omega|^2/2}$ at time $1$, run at the
random time of the subordinator with exponent $u \mapsto f(2u)$, has transform $e^{-f(|\omega|^2)}$
at time $1$. Its statement is typed in the skeleton (\texttt{PyramidsSkeleton.Ch02CitedFacts}) and is not an axiom; admitting it is the author's decision at phase~5. \textbf{Assignment.}
\ledger{A57} carries the statement. No node spends it: chapter~4 states it in the paragraph before
Proposition~\ref{prop:matern-increments}, whose clause (2) writes the laws down directly. Not
carried: the converse of the theorem, that negative definiteness in every dimension makes $f$ a
Bernstein function; and the killed case $f(0{+}) > 0$.)}
\end{theorem}

\subsection{\Polya{} frequency densities}
\label{subsec:cited-karlin}

% draft: the citation of [Karlin 1968, Ch. 7, Thm. 3.2(a), p. 345, with Prop. 1.4, p. 333] at line
% 235, as printed in [V-selection v0.1, section 4.3], and the same fact at line 255.
% CHANGED: the ranges gamma >= 0, delta real, a_j real and Karlin's normalization
% gamma + sum a_j^2 > 0, which the draft omits, are stated; "analytic on the strip" is the page's
% "converges absolutely on the strip" (VERBATIM-selection-iso-affine.md and
% VERBATIM-literature.md, A24).
% CHANGED (fidelity review, 2026-11-15, R52): the condition gamma + sum a_j^2 > 0 is anchored to the
% statement of Thm. 3.2(a), p. 345, where it stands; Remark 3.1, p. 346, to which V-selection p. 35
% attributes it, is about dropping integrability. The statement is unchanged.
\begin{theorem}[Schoenberg's representation]
\label{thm:cited-A24}
\lean{Pyramids.cited_A24}
\leanok
% CHANGED (2026-11-15, admitted narrower than the page, ADR-0005): the identity is stated on the
% imaginary axis, where both spenders read it. It was: "Then there are an open strip
% alpha < Re s < beta with alpha < 0 < beta, on which L converges absolutely, and numbers ...,
% such that on the strip L(s) psi(s) = 1, ...".
Let $\phi$ be a \Polya{} frequency density on $\RR$, and $L(s) := \int_{\RR}e^{-sx}\phi(x)\,dx$ its
two-sided Laplace transform. Then there are numbers $\gamma \ge 0$, $\delta \in \RR$ and real $a_j$
with $\sum_j a_j^2 < \infty$ and $\gamma + \sum_j a_j^2 > 0$ such that
\[
  L(i\omega)\,\psi(i\omega) = 1 \quad (\omega \in \RR), \qquad
  \psi(s) = e^{-\gamma s^2 + \delta s}\prod_j(1 + a_j s)\,e^{-a_j s},
\]
where $\psi$ is an entire function with $\psi(0) = 1$.
% CHANGED (fidelity review, 2026-11-15, R26/R50/R69): the status line lists the declarations' hypotheses beside the node's (LINKAGE rule 6).
% CHANGED (fidelity review, 2026-11-16, R57): the PF hypothesis is "same", not "stronger": Karlin's
% definition (Ch. 1, pp. 11-12, rendered and read) requires the determinants at every ordered tuple.
\statusA\quad\ledger{A24}\quad\emph{(Cited from [\sscite{karlin1968total}, Ch.~7, Thm.~3.2(a),
p.~345], whose statement carries the condition $\gamma + \sum_ja_j^2 > 0$, and the class
$\mathcal{E}_2^*$ of its display (2.2), p.~336; as printed in
\prog{V-selection v0.1}{\S~4.3, pp.~34--36}, which reads it at Karlin's letter, in Karlin's sign of
the transform, and anchors the condition to Remark~3.1, p.~346, which is about dropping
integrability. The page gives the identity $L\psi = 1$ on an open strip
$\alpha < \mathrm{Re}\,s < \beta$ containing the imaginary axis, on which $L$ exists
(Prop.~1.4, p.~333); and only the identity there is stated and
admitted. Admitted as a Lean axiom by the author on 2026-11-15, at this statement.
\textbf{Assignment.} \ledger{A24} carries the statement.
\ref{thm:matern-pyramid-line-noncreation} spends that the reciprocal of $L$ on the imaginary axis
is the restriction of the entire function $\psi$, and
\ref{prop:classical-pyramids-approximate-discrete-gaussian} the modulus of $\psi$ there. Not
carried: the strip and the absolute convergence on it; the converse, that $1/\psi$ is the Laplace transform of a \Polya{}
frequency density for every such $\psi$; and the claim that a \Polya{} frequency density has an
absolutely continuous \Levy{} measure, which the selection module does not print (plan id A25) and
\ref{prop:classical-pyramids-approximate-discrete-gaussian} proves.
\textbf{Hypotheses} (the node's, each with what its declarations require): $\phi$ a \Polya{} frequency
density on $\RR$: integrable with integral $1$, same; \Polya{} frequency, same,
\texttt{IsPolyaFrequencyDensity} requiring $\det[\phi(x_i - y_j)] \ge 0$ at every pair of increasing
point families, of every order, as Karlin's definition does (\sscite{karlin1968total}, Ch.~1,
\S1, pp.~11--12: $K(x, y)$ is $\mathrm{TP}_r$ if the determinants (1.2) are nonnegative ``for
all'' increasing $x_1 < \dots < x_m$, $y_1 < \dots < y_m$, and $f$ is $\mathrm{PF}_r$ if
$f(x - y)$ is $\mathrm{TP}_r$); every spender obtains $\phi$ from \texttt{cited\_A22\_polya},
which gives that form.)}
\end{theorem}

\subsection{The modified Bessel function}
\label{subsec:cited-bessel}

The modified Bessel function $K_\nu$ is taken to be defined by the integral of
[\sscite{dlmf2026}, 10.32.9], so the next theorem says that this integral equals the one of
[\sscite{dlmf2026}, 10.32.10]. The source states the second representation for complex $z$ with
$|\mathrm{ph}\,z| < \pi/4$; the theorem takes real $z > 0$.

% draft: the citation of [DLMF 10.32.10] at line 541; the evaluations at lines 590 and 602 are
% proved in lem:bessel-k from 10.32.10 and the definition 10.32.9.
% CHANGED (2026-11-13, the author's decision): the former clause (2), the further DLMF equations
% 10.29.3, 10.32.8, 10.39.2, 10.29.1 and 10.30.3 with their Lean axiom cited_A51_evaluations, is
% dropped, as no node and no proof read it (a safe-direction change; the ledger keeps a row).
\begin{theorem}[the modified Bessel function $K_\nu$]
\label{thm:cited-A51}
\lean{Pyramids.cited_A51_integral}
\leanok
% CHANGED (fidelity review, 2026-11-16, R75): the body says that $K_\nu$ is taken as the integral of DLMF 10.32.9, so that the node's content is that 10.32.9 equals 10.32.10.
% CHANGED (the theory paper's blueprint round, 2026-11-16, review of Section 2, ENVELOPE): stated
% directly, K_nu defined by 10.32.9 and the equation 10.32.10; the sentence on how the statement is
% read moves to the prose before the environment. Same content.
% CHANGED (skeleton review, M2): narrowed to real z > 0, the range the one spender reads; the
% page's complex sector |ph z| < pi/4 is kept in the annotation as what it prints.
For real $\nu$ and real $z > 0$ let
$K_\nu(z) := \int_0^\infty e^{-z\cosh u}\cosh(\nu u)\,du$, the modified Bessel function of the
second kind of order $\nu$ (DLMF~10.32.9). Then for every real $\nu$ and every real $z > 0$
\[
  K_\nu(z) = \tfrac12\bigl(\tfrac12 z\bigr)^{\nu}\int_0^\infty
  \exp\Bigl(-t - \frac{z^2}{4t}\Bigr)\,\frac{dt}{t^{\nu+1}}
\]
(DLMF~10.32.10).
% CHANGED (fidelity review, 2026-11-16, R71): the status line lists the declaration's hypotheses beside the node's (LINKAGE rule 6).
\statusA\quad\ledger{A51}\quad\emph{(Cited from [\sscite{dlmf2026}], the NIST Digital Library of
Mathematical Functions, Version 1.2.8 of 2026-09-15, read online: equation 10.32.10,
https://dlmf.nist.gov/10.32.E10, which prints no condition on $\nu$ and holds there for complex
$z$ with $|\mathrm{ph}\,z| < \tfrac14\pi$, of which the statement takes the real half-line.
Admitted by the author on 2026-11-13. The typed $K_\nu$ is the shared core's, defined at real
$z > 0$ by the integral $\int_0^\infty e^{-z\cosh u}\cosh(\nu u)\,du$, which is DLMF~10.32.9 on the
real half-line (its wording is transcribed in the ledger). The statement is proved in Lean from
that definition, by the substitution $u = c e^{v}$ and a Gaussian bound that makes the integral
finite, and needs no axiom.
\textbf{Assignment.} \ledger{A51} carries the formula. \ref{lem:bessel-k} spends it at real
$z = y > 0$ and $\nu = 0$, $1$ and $d/2$. Not carried: anything about $K_\nu$ at complex argument,
and the further DLMF equations the draft reads at its lines 590 and 602, which
\ref{lem:bessel-k} proves.
\textbf{Hypotheses} (the node's, each with what its declaration requires): $\nu$ real: same, every real
$\nu$; $z$ real with $z > 0$: same (at $z \le 0$ the definition of $K_\nu$ returns the junk value
$0$, which is why the hypothesis is required).)}
\end{theorem}
